
%%  FINAL YEAR PROJECT REPORT
%%  Design and Implementation of an IoT-Enabled Inverter System
%%  with Battery-Aware Load Prioritization and Remote Monitoring
%%
%%  Kayode Daniel Oluwatobi | EEE/20/5675
%%  Federal University of Technology, Akure (FUTA)
%%  Department of Electrical and Electronics Engineering
%%  AUGUST, 2026
%% ============================================================

\documentclass[12pt, a4paper]{report}

%% ── Packages ─────────────────────────────────────────────────
\usepackage[
  a4paper,
  left=1.1in,
  right=1.0in,
  top=1.0in,
  bottom=1.0in,
  footskip=0.6in
]{geometry}

\usepackage{times}
\usepackage{setspace}
\usepackage{graphicx}
\usepackage{amsmath}
\usepackage{amssymb}
\usepackage{array}
\usepackage{caption}
\usepackage{subcaption}
\usepackage{float}
\usepackage[hidelinks]{hyperref}
\usepackage{url}
\usepackage{enumitem}
\usepackage{titlesec}
\usepackage{fancyhdr}
\usepackage{listings}
\usepackage{xcolor}
\usepackage{microtype}
\usepackage{tabularx}
\usepackage{booktabs}
\usepackage{tikz}
\usepackage{tocloft}
\usepackage{emptypage}
\usepackage{longtable}
\usepackage[authoryear, round]{natbib}

\usetikzlibrary{shapes.geometric, arrows.meta, positioning, fit, backgrounds}

%% ── Hyphenation prevention ───────────────────────────────────
\hyphenpenalty=10000
\exhyphenpenalty=10000
\tolerance=1000
\emergencystretch=3em

%% ── Widow and orphan control ─────────────────────────────────
\widowpenalty=10000
\clubpenalty=10000
\displaywidowpenalty=10000
\raggedbottom

%% ── Line spacing ─────────────────────────────────────────────
\onehalfspacing

%% ── Block paragraphing ───────────────────────────────────────
\setlength{\parindent}{0pt}
\setlength{\parskip}{10pt}



% Reduce vertical spacing between chapter-level TOC entries
\setlength{\cftbeforechapskip}{0pt}
%% Reduce space above TOC title
\setlength{\cftaftertoctitleskip}{30pt}

%% ── Chapter heading format ───────────────────────────────────
\titleformat{\chapter}[display]
  {\normalfont\fontsize{14}{16}\bfseries\centering}
  {\MakeUppercase{\chaptertitlename}\ \thechapter}
  {-5pt}
  {\fontsize{14}{16}\bfseries\MakeUppercase}

\titleformat{\section}
  {\normalfont\fontsize{13}{15}\bfseries\raggedright}
  {\thesection}
  {0.6em}
  {\MakeUppercase}

\titleformat{\subsection}
  {\normalfont\fontsize{12}{14}\bfseries\raggedright}
  {\thesubsection}
  {0.6em}
  {}

\titleformat{\subsubsection}
  {\normalfont\fontsize{12}{14}\bfseries\itshape\raggedright}
  {\thesubsubsection}
  {0.6em}
  {}

\titlespacing*{\chapter}{0pt}{-20pt}{10pt}
\titlespacing*{\section}{0pt}{16pt}{8pt}
\titlespacing*{\subsection}{0pt}{12pt}{6pt}
\titlespacing*{\subsubsection}{0pt}{10pt}{5pt}

%% ── Table of contents format ─────────────────────────────────
\renewcommand{\contentsname}{TABLE OF CONTENTS}
\setlength{\cftbeforetoctitleskip}{-20pt}
\renewcommand{\cfttoctitlefont}{\hfill\fontsize{14}{16}\bfseries\MakeUppercase}
\renewcommand{\cftaftertoctitle}{\hfill}
\renewcommand{\cftloftitlefont}{\hfill\fontsize{14}{16}\bfseries\MakeUppercase}
\renewcommand{\cftafterloftitle}{\hfill}
\setlength{\cftbeforeloftitleskip}{-20pt}
\renewcommand{\cftlottitlefont}{\hfill\fontsize{14}{16}\bfseries\MakeUppercase}
\renewcommand{\cftafterlottitle}{\hfill}
\setlength{\cftbeforelottitleskip}{-20pt}

%% No bold in TOC
\renewcommand{\cftchapfont}{\normalfont}
\renewcommand{\cftchappagefont}{\normalfont}
\renewcommand{\cftsecfont}{\normalfont}
\renewcommand{\cftsecpagefont}{\normalfont}
\renewcommand{\cftsubsecfont}{\normalfont}
\renewcommand{\cftsubsecpagefont}{\normalfont}

%% No dotted leaders in TOC
\renewcommand{\cftchapleader}{\hfill}
\renewcommand{\cftsecleader}{\hfill}
\renewcommand{\cftsubsecleader}{\hfill}
\renewcommand{\cftsubsubsecleader}{\hfill}

%% No dotted leaders in List of Figures
\renewcommand{\cftfigleader}{\hfill}

%% No dotted leaders in List of Tables
\renewcommand{\cfttableader}{\hfill}

%% Chapter prefix in TOC: CHAPTER 1: TITLE
\renewcommand{\cftchappresnum}{CHAPTER }
\renewcommand{\cftchapaftersnum}{:}
\setlength{\cftchapnumwidth}{6.5em}

%% ── Page numbering ───────────────────────────────────────────
\pagestyle{fancy}
\fancyhf{}
\cfoot{\thepage}
\renewcommand{\headrulewidth}{0pt}
\fancypagestyle{plain}{
  \fancyhf{}
  \cfoot{\thepage}
  \renewcommand{\headrulewidth}{0pt}
}

%% ── Colours ──────────────────────────────────────────────────
\definecolor{codebg}{RGB}{248,248,248}
\definecolor{codeframe}{RGB}{210,210,210}
\definecolor{codecomment}{RGB}{106,106,106}
\definecolor{codekeyword}{RGB}{0,0,160}
\definecolor{codestring}{RGB}{163,21,21}
\definecolor{blockblue}{RGB}{220,232,255}
\definecolor{blockgreen}{RGB}{210,240,210}
\definecolor{blockorange}{RGB}{255,237,210}
\definecolor{blockgrey}{RGB}{235,235,235}
\definecolor{arrowgrey}{RGB}{90,90,90}

%% ── Code style ───────────────────────────────────────────────
\lstdefinestyle{mycode}{
  backgroundcolor=\color{codebg},
  commentstyle=\color{codecomment}\itshape,
  keywordstyle=\color{codekeyword}\bfseries,
  stringstyle=\color{codestring},
  basicstyle=\ttfamily\small,
  breakatwhitespace=false,
  breaklines=true,
  captionpos=b,
  keepspaces=true,
  numbers=left,
  numbersep=6pt,
  numberstyle=\tiny\color{gray},
  showstringspaces=false,
  tabsize=2,
  frame=single,
  rulecolor=\color{codeframe},
  xleftmargin=12pt,
  xrightmargin=4pt,
}


\lstset{style=mycode}


%% Remove dotted leaders from List of Listings
\makeatletter
\renewcommand*\l@lstlisting{\@dottedtocline{1}{1.5em}{3.2em}}
\renewcommand{\@dotsep}{500}  % effectively removes dots
\makeatother

\lstdefinelanguage{JavaScript}{
  keywords={const,let,var,function,return,if,else,for,while,async,
             await,import,export,default,class,new,this,true,false,null},
  morecomment=[l]{//},
  morecomment=[s]{/*}{*/},
  morestring=[b]',
  morestring=[b]",
  sensitive=true
}

\lstdefinelanguage{CppESP}{
  language=C++,
  morekeywords={Serial,WiFi,HTTPClient,String,float,int,bool,
                uint32_t,uint8_t,void,setup,loop,digitalWrite,
                analogRead,delay,millis},
  sensitive=true
}

%% ── TikZ styles ──────────────────────────────────────────────
\tikzset{
  sysblock/.style={
    rectangle, rounded corners=4pt,
    minimum width=2.8cm, minimum height=0.8cm,
    text centered, font=\small\bfseries,
    draw=black!60, line width=0.6pt
  },
  arrow/.style={-Stealth, thick, color=arrowgrey},
  dasharrow/.style={-Stealth, thick, dashed, color=arrowgrey}
}

%% ── Table style: no vertical lines, top+bottom+header rules ─
\newcolumntype{L}[1]{>{\raggedright\arraybackslash}p{#1}}
\newcolumntype{C}[1]{>{\centering\arraybackslash}p{#1}}

%% ── Hyperref ─────────────────────────────────────────────────
\hypersetup{colorlinks=true,linkcolor=black,citecolor=black,urlcolor=blue}

%% ─────────────────────────────────────────────────────────────
\begin{document}

%% ══ TITLE PAGE (page i, invisible) ══════════════════════════
\pagenumbering{roman}
\thispagestyle{empty}
\begin{center}
  \vspace*{0.1cm}
  \includegraphics[width=4.5cm]{futalogo.jpeg}

  \vspace{0.8cm}

  {\fontsize{14}{20}\bfseries
    DESIGN AND IMPLEMENTATION OF AN IoT-ENABLED INVERTER SYSTEM
    WITH BATTERY-AWARE LOAD PRIORITIZATION
    AND REMOTE MONITORING
  }

  \vspace{0.9cm}

  {\normalsize BY}

  \vspace{0.1cm}

  {\fontsize{12}{16} DANIEL OLUWATOBI, KAYODE }\\[0.3em]
  {\fontsize{12}{16} EEE/20/5675}

  \vspace{0.8cm}

  {\normalsize SUPERVISOR:}\\[0.2em]
  {\normalsize DR. O. A. AGBOLADE}

  \vspace{0.9cm}

  {\normalsize SUBMITTED TO}\\[0.2em]
  
 \vspace{0.4cm}
 
  {\normalsize
    THE DEPARTMENT OF ELECTRICAL AND ELECTRONICS ENGINEERING,\\[0.2em]
    SCHOOL OF ENGINEERING AND ENGINEERING TECHNOLOGY,\\[0.2em]
    FEDERAL UNIVERSITY OF TECHNOLOGY, AKURE.
  }

  \vspace{0.6cm}

  {\normalsize
    IN PARTIAL FULFILLMENT OF THE REQUIREMENTS FOR THE AWARD OF\\[0.2em]
    A BACHELOR OF ENGINEERING (B.ENG) DEGREE IN\\[0.2em]
    ELECTRICAL AND ELECTRONICS ENGINEERING.
  }

 \vfill
\begin{flushright}
  {\normalsize AUGUST, 2026}
\end{flushright}
\end{center}

%% ══ CERTIFICATION (page ii) ═════════════════════════════════
\newpage
\setcounter{page}{2}
\thispagestyle{plain}
\addcontentsline{toc}{chapter}{CERTIFICATION}

\begin{center}
  {\fontsize{14}{16}\bfseries CERTIFICATION}
\end{center}

\vspace{10pt}

This document certifies that this report provides a comprehensive record
on the final-year undergraduate project work done by KAYODE DANIEL
OLUWATOBI, matriculation number EEE/20/5675, in the Department of
Electrical and Electronics Engineering, the Federal University of
Technology, Akure. The report has been prepared following the regulations
governing the production of reports in the department and submitted in
partial fulfillment of the requirement for the award of a Bachelor of
Engineering (B.Eng.) degree in Electrical and Electronics Engineering.

\vspace{60pt}

\noindent
\begin{minipage}[t]{0.45\textwidth}
  Kayode Daniel Oluwatobi\\[2pt]
  (Student)
\end{minipage}
\hfill
\begin{minipage}[t]{0.45\textwidth}
  \dotfill\\[2pt]
  \centering Signature and Date
\end{minipage}

\vspace{90pt}

\noindent
\begin{minipage}[t]{0.45\textwidth}
  Dr. O. A. Agbolade\\[2pt]
  (Supervisor)
\end{minipage}
\hfill
\begin{minipage}[t]{0.45\textwidth}
  \dotfill\\[2pt]
  \centering Signature and Date
\end{minipage}

\vspace{90pt}

\noindent
\begin{minipage}[t]{0.45\textwidth}
  Dr. B. C. Ubochi\\[2pt]
  (Head of Department)
\end{minipage}
\hfill
\begin{minipage}[t]{0.45\textwidth}
  \dotfill\\[2pt]
  \centering Signature and Date
\end{minipage}

%% ══ DEDICATION (page iii) ═══════════════════════════════════
\newpage
\setcounter{page}{3}
\thispagestyle{plain}
\addcontentsline{toc}{chapter}{DEDICATION}

\begin{center}
  {\fontsize{14}{16}\bfseries DEDICATION}
\end{center}

\vspace{10pt}

This work is dedicated to  the Almighty God, to my beloved parents, Dr. and Mrs Kayode, and my supervisor, Dr. O. A. Agbolade

%% ══ ACKNOWLEDGEMENT (page iv) ══════════════════════════════
\newpage
\setcounter{page}{4}
\thispagestyle{plain}
\addcontentsline{toc}{chapter}{ACKNOWLEDGEMENT}

\begin{center}
  {\fontsize{14}{16}\bfseries ACKNOWLEDGEMENT}
\end{center}

\vspace{10pt}

All glory and honour belong to Almighty God, the source of all knowledge
and wisdom, without whose grace this project and the entire undergraduate
programme would not have been possible. I am deeply grateful for His
provision, direction and strength throughout every phase of this work.

I owe an immense debt of gratitude to my parents, who believed in me from
the very beginning. Their phone calls to check in on my progress, their
words of encouragement during difficult moments and the sacrifices they
made to fund this education are the foundation upon which this
achievement rests. Thank you for everything.

I am sincerely grateful to my supervisor, Dr. O. A. Agbolade, for his
guidance, patience and dedication at every step of this project. His
technical advice, constructive feedback and willingness to be available
whenever the team needed direction made a profound difference to the
quality of this work. I could not have asked for a more committed
supervisor.

I also acknowledge the Head of Department, Dr. B. C. Ubochi, and the
entire academic and non-academic staff of the Department of Electrical
and Electronics Engineering at the Federal University of Technology,
Akure. Their collective efforts throughout my five years of undergraduate
study built the foundation of knowledge and practical skill that this
project draws upon. I am grateful for every lecture, laboratory session
and administrative support that contributed to my development as an
engineer.

%% ══ ABSTRACT (page v) ══════════════════════════════════════
\newpage
\setcounter{page}{5}
\thispagestyle{plain}
\addcontentsline{toc}{chapter}{ABSTRACT}

\begin{center}
  {\fontsize{14}{16}\bfseries ABSTRACT}
\end{center}

\vspace{10pt}

Incessant power grid instability in developing countries has necessitated widespread adoption of solar-inverter backup systems. However, conventional installations frequently suffer from premature battery bank degradation due to unmonitored over-discharge, lack of cell-level observability, and indiscriminate supplying of non-critical loads during outages. This project presents the design, implementation, and empirical validation of an Internet of Things (IoT)-enabled smart inverter management system featuring real-time cell telemetry, automated battery-aware load prioritization, and cloud-synchronized remote monitoring and control.

The hardware layer comprises a 4-series 100\,Ah Lithium Iron Phosphate ($\text{LiFePO}_4$) prismatic battery pack ($12.8\,\text{V}$, $1.28\,\text{kWh}$) mechanically constrained with custom steel compression plates and insulating polymer spacers, managed by a smart JK Battery Management System (BMS) with active inductive cell balancing. To circumvent Bluetooth and Wi-Fi radio co-existence conflicts on a single microcontroller, a dual-ESP32 architecture was designed: ESP32 Board~1 maintains a low-latency NimBLE connection to the BMS to reassemble 300-byte fragmented telemetry frames and drive a local 2.4-inch ILI9341 SPI TFT display, while ESP32 Board~2 handles Wi-Fi connectivity, multi-channel active-HIGH relay switching, and analog current measurement via multi-sampled ACS758 Hall-effect sensors on dedicated ADC1 channels with boot-time zero-point calibration. A 3-stage load prioritization engine autonomously sheds non-essential ($60\%$), major ($40\%$), and critical ($10\%$) load channels with a $5\%$ hysteresis buffer to protect pack longevity. The software ecosystem features bi-directional UART communication, cloud publishing to Firebase Realtime Database for live states and Supabase for historical analytics, and an installable Next.js 14 Progressive Web Application with Google OAuth and Web Speech voice alarms.

Experimental validation demonstrated $100\%$ packet reassembly reliability over BLE, flawless zero-point sensor calibration ($<0.1\,\text{A}$ noise deadband), and precise automated load shedding across all State of Charge (SoC) thresholds under simulated load conditions. Cell voltage differentials remained within $3\,\text{mV}$ to $17\,\text{mV}$ under active balancing, confirming cell uniformity. Remote command dispatch and telemetry synchronization achieved end-to-end latencies under 3~seconds and 5~seconds, respectively. The system successfully demonstrates a robust, cost-effective architecture for extending energy storage lifespan and optimizing backup runtime in residential and commercial energy systems.

%% ══ TABLE OF CONTENTS ══════════════════════════════════════
\newpage
\setcounter{page}{6}
\addcontentsline{toc}{chapter}{TABLE OF CONTENTS}
\tableofcontents

%% ══ LIST OF FIGURES ═════════════════════════════════════════
\newpage
\listoffigures

%% ══ LIST OF TABLES ══════════════════════════════════════════
\newpage
\listoftables
%% ══ LIST OF LISTINGS ════════════════════════════════════════
\newpage
\renewcommand{\lstlistlistingname}{LIST OF LISTINGS}
\lstlistoflistings


%% ══════════════════════════════════════════════════════════════
%%  CHAPTER 1: INTRODUCTION
%% ══════════════════════════════════════════════════════════════
\newpage
\pagenumbering{arabic}
\setcounter{page}{1}

\chapter{INTRODUCTION}

\section{Background of the Study}

Electricity is the backbone of economic activity and quality of life in
the modern world. In Nigeria, however, the reliability of the national
grid has long been a persistent and well-documented challenge. The country
operates one of the largest but least reliable power grids in sub-Saharan
Africa, with per-capita electricity consumption among the lowest in the
region despite significant installed generation capacity \citep{aliyu2013}.
Frequent outages, voltage instability and load shedding events impose
enormous costs on households, small businesses and industrial operators
alike.

This energy deficit has catalysed widespread adoption of alternative power
backup systems. Among these, battery-backed inverter systems have grown
substantially in popularity over petrol-powered generators because they
operate silently, produce no exhaust fumes and can be recharged from the
grid or from solar photovoltaic panels whenever supply is available
\citep{rashid2014}. The inverter converts stored Direct Current energy from
a battery bank into Alternating Current suitable for household appliances
and office equipment.

Despite their utility, conventional inverter systems share a fundamental
limitation: once installed, they operate as passive devices offering the
user very little visibility or control. A user who steps out of the house
has no way to determine how much battery capacity remains, which loads are
drawing power or whether the battery is approaching a dangerous discharge
level. When the battery is eventually depleted, the system shuts down
without warning. Repeated deep discharge events of this nature permanently
damage lithium battery cells, reducing capacity and shortening the
lifespan of what is typically an expensive investment \citep{omar2014}.

The Internet of Things represents a significant technological shift that
directly addresses these limitations. IoT refers to the embedding of
communication hardware in physical devices so that they can exchange data
over wireless networks, enabling remote monitoring, automated decision
making and real-time control from anywhere in the world \citep{ashton2009}.
When applied to a home inverter system, IoT technology makes it possible
to monitor battery State of Charge in real time, receive push alerts
before a critical condition develops, automatically manage which loads
remain powered as the battery depletes and issue control commands remotely
from a smartphone.

This project implements exactly that vision. It presents the design and
implementation of a complete IoT-enabled inverter system that integrates a
Lithium Iron Phosphate battery pack, a Battery Management System, a dual
microcontroller embedded firmware architecture and a cloud-connected
Progressive Web Application, all working in concert to deliver
intelligent, remotely accessible power management to a Nigerian household.

\section{Problem Statement}

Conventional battery-backed inverter systems deployed in Nigerian homes
operate without any real-time monitoring or remote control capability.
Users have no visibility into battery State of Charge, instantaneous power
consumption or individual load behaviour except by physically reading a
display at the inverter unit. When power is required remotely or during
extended absences, the user is entirely unable to assess the system's
status or intervene to protect the battery from damage through over-discharge.

Furthermore, most systems apply load shedding uniformly: when the battery
depletes, the entire output is disconnected, cutting power to critical
devices such as routers, medical equipment and lighting simultaneously
with non-essential loads such as air conditioners and entertainment
systems. This indiscriminate shutdown reduces the effective usefulness of
the backup system and accelerates battery degradation.

There is also no mechanism for historical analysis of energy consumption
patterns, which prevents users from making informed decisions about load
scheduling, battery capacity requirements or tariff optimisation.

\section{Aim of the Project}

This project aims to develop an Internet of Things (IoT)-enabled inverter system for real-time battery monitoring, automated load shedding, and remote cloud control, improving battery lifespan and user experience.


\section{Objectives of the Project}

The objectives of this project are to:

\begin{enumerate}[noitemsep, leftmargin=0.7in, label=\roman*.]
  \item design and assemble a 4-series Lithium Iron Phosphate battery
        pack and integrate the JK Battery Management System for cell-level
        monitoring and protection.
  \item develop firmware for two ESP32 microcontrollers implementing
        BLE communication with the JK BMS, real-time data parsing,
        UART inter-board communication, Wi-Fi connectivity, relay-based
        load switching and cloud telemetry publishing.
  \item implement an intelligent load shedding engine that
        automatically disconnects output channels in priority order as
        battery State of Charge decreases below configurable thresholds.
  \item develop a cloud-connected Progressive Web Application that
        displays live and historical battery telemetry, provides remote
        load channel control and supports MOSFET command dispatch through
        Firebase Realtime Database.
\end{enumerate}

\section{Scope of the Project}

This project covers the complete design and implementation of an
IoT-enabled inverter monitoring system, from battery pack assembly through
embedded firmware development to cloud infrastructure and web application
deployment. The hardware components include a 4S LiFePO\textsubscript{4}
battery bank, the JK BMS, two ESP32 development boards, a 2.4-inch TFT
display, three ACS758 Hall-effect current sensors, three relay modules and
the inverter output wiring.

The firmware covers BLE communication and binary frame parsing for the
JK BMS, TFT display rendering, UART inter-board communication, load
prioritization control, ACS758 current measurement, Firebase and Supabase
cloud publishing and NTP time synchronisation. The web application covers
real-time dashboard display, historical analytics charts and remote relay
and MOSFET control.

The project does not cover the internal design of the inverter power stage,
the design of a custom Battery Management System or the implementation of
solar charge control. These elements exist in the system but use
commercially available units whose internal design is outside the scope
of this report.

\section{Significance of the Project}

This project makes several contributions to the field of intelligent
embedded energy systems. It demonstrates that a complete IoT-enabled
battery monitoring and load management system can be realised using
accessible, low-cost hardware components within the resource constraints
of an undergraduate final-year project in Nigeria.

From a technical standpoint, the dual-ESP32 architecture resolves a
real hardware limitation: the inability of a single ESP32 to run
Bluetooth Low Energy and Wi-Fi simultaneously without severe packet loss.
The firmware-level implementation of the JK BMS proprietary binary
protocol provides a working reference for future researchers seeking to
interface with JK BMS devices programmatically. The use of ADC1-only
pins for current sensing while Wi-Fi is active addresses a commonly
overlooked hardware conflict in ESP32-based IoT designs.

From a practical standpoint, the system directly addresses a pressing
need in the Nigerian context: giving households real-time visibility and
remote control over their inverter systems, extending battery lifespan
through protected load management and enabling data-driven energy
consumption analysis.

%% ══════════════════════════════════════════════════════════════
%%  CHAPTER 2: LITERATURE REVIEW
%% ══════════════════════════════════════════════════════════════
\newpage
\chapter{LITERATURE REVIEW}

\section{Lithium Iron Phosphate Battery Technology}

Lithium Iron Phosphate (LiFePO\textsubscript{4}) batteries belong to the
broader family of lithium-ion electrochemical cells and have emerged as
the preferred chemistry for stationary energy storage, backup power systems
and electric vehicle applications where safety, longevity and thermal
stability are priorities. The chemistry uses lithium iron phosphate as the
cathode material, graphite as the anode and a lithium salt dissolved in an
organic solvent as the electrolyte \citep{nitta2015}.

Compared to other lithium-ion chemistries such as Lithium Cobalt Oxide and
Nickel Manganese Cobalt, LiFePO\textsubscript{4} offers superior thermal
and chemical stability. The iron-phosphate bond in the cathode structure is
significantly stronger than the bonds in competing chemistries, which means
the material does not decompose at elevated temperatures and is far less
susceptible to thermal runaway \citep{dini2024}. This makes it
substantially safer for unattended residential installations, which is
the primary deployment context of this project.

LiFePO\textsubscript{4} cells nominally operate at 3.2V per cell, reaching
approximately 3.65V at full charge and dropping to a safe lower cutoff of
about 2.5V per cell \citep{piller2001}. In a four-cell series configuration
as used in this project, these limits translate to a nominal pack voltage
of 12.8V, a full-charge voltage of 14.6V and a lower cutoff of 10V.
The discharge curve of LiFePO\textsubscript{4} is notably flat across the
bulk of the capacity range, remaining near 3.2V to 3.3V per cell between
approximately 20\% and 80\% State of Charge. This flat profile
complicates voltage-based State of Charge estimation and is one of the
reasons that the JK BMS relies on coulomb counting as its primary
estimation method, supplemented by voltage cross-checks.

\citet{pillet2025} investigated the long-term cycling behaviour of
LiFePO\textsubscript{4} cells under varying temperature and charge rate
conditions and confirmed that the chemistry delivers excellent cycle life
when operated within its recommended voltage window, with capacity
retention above 80\% after 2000 cycles under nominal conditions.

\subsection{Mechanical Compression of Prismatic LiFePO\texorpdfstring{\textsubscript{4}}{4} Cells}

Prismatic LiFePO\textsubscript{4} cells are housed in thin aluminium
enclosures. During charge and discharge cycles, the intercalation and
de-intercalation of lithium ions cause the electrode materials to expand
and contract, producing a reversible volumetric change in the cell.
Without mechanical constraint, this repeated dimensional cycling causes
progressive delamination of the electrode layers from each other and from
their current collectors \citep{fraunhofer2022}. Delamination increases
internal resistance, reduces usable capacity and, in severe cases, causes
permanent structural collapse of the electrode assembly.

Mechanical compression using rigid steel end plates and threaded steel
rods restrains this volumetric change, keeping the electrode layers in
intimate contact and preventing the irreversible delamination that would
otherwise shorten the cell's useful life. The compression fixture applies
a controlled and uniform pressure across the cell faces, distributing the
internal stress evenly rather than allowing localised swelling to
accumulate. Research by \citet{kit2024} demonstrated that appropriate
compression levels extend cycle life significantly compared to uncompressed
cells, while excessive compression introduces its own stress-related
degradation pathway. In this project, steel compression plates and
threaded rods are used to compress the four prismatic cells together,
maintaining structural integrity throughout the charging and discharging
cycles.

\section{Battery Management Systems}

A Battery Management System is an electronic circuit assembly that
monitors and protects a lithium battery pack during operation. The
functions of a BMS include cell voltage monitoring, pack current
measurement, temperature monitoring, State of Charge estimation,
overcharge protection, over-discharge protection, overcurrent protection,
short-circuit protection and cell balancing \citep{dini2024}. Without a
BMS, lithium cells are vulnerable to all of these failure modes, each of
which can permanently damage the cells or create a fire hazard.

\citet{hannan2017} provided a comprehensive review of BMS technologies and
energy management strategies for lithium-ion battery packs, highlighting
that accurate State of Charge estimation remains one of the most critical
and technically challenging functions. Several estimation methods have
been developed and deployed. Coulomb counting integrates the measured
current over time to track the flow of charge in and out of the battery.
It is simple and computationally inexpensive but accumulates errors over
long periods due to current sensor inaccuracies and self-discharge effects.
The Open-Circuit Voltage method estimates State of Charge from the
battery's resting terminal voltage using a pre-characterised voltage versus
State of Charge lookup table, but requires the battery to be at rest for
an extended period for the measurement to be valid. Model-based methods
such as Extended Kalman Filters combine current integration with voltage
correction in a probabilistic framework that is more accurate but
significantly more complex to implement \citep{piller2001}. The JK BMS
used in this project employs coulomb counting as its primary State of
Charge estimator, periodically recalibrated against the terminal voltage
when the battery is at a known charge state.

\subsection{Active Cell Balancing}

Cell balancing addresses the practical reality that individual cells in a
series-connected pack are never perfectly identical. Manufacturing
tolerances produce small differences in capacity, internal resistance and
self-discharge rate that cause the cells to drift apart in State of Charge
over repeated cycles. Because the BMS must protect the weakest cell from
exceeding its voltage limits, the entire pack's usable capacity is
constrained by that weakest cell. Passive balancing dissipates excess
charge from higher-voltage cells as heat through a resistor, equalising
cell voltages at the cost of wasted energy. Active balancing, as
implemented in the JK BMS, transfers charge from higher-voltage cells to
lower-voltage cells using an inductor-based energy transfer circuit,
achieving equalisation without wasting the energy. \citet{dini2024}
demonstrated that active balancing substantially outperforms passive
balancing in both energy efficiency and long-term capacity retention.

\section{ESP32 Microcontroller in IoT Applications}

The ESP32 is a dual-core 32-bit microcontroller developed by Espressif
Systems, integrating both Wi-Fi (802.11 b/g/n) and Bluetooth (Classic and
BLE 4.2) radios in a single chip \citep{espressif2024}. This integration
makes it particularly attractive for IoT applications where wireless
connectivity is essential but board space and cost constraints preclude
separate radio modules.

\citet{dzahir2023} evaluated the energy consumption of ESP32
microcontrollers in real-time MQTT-based IoT monitoring systems under
four operating conditions: always-on, light sleep, deep sleep with
periodic wake and modem sleep. Their study found that deep sleep mode
extends battery runtime from a few days to over 17 days, confirming the
importance of power management strategy in deployed IoT nodes. While
the present project does not prioritise power minimisation of the ESP32
boards (they are powered from the inverter system itself), these findings
inform the rationale for separating BLE and Wi-Fi radio operations across
two boards, as the radio activity is itself a major contributor to power
consumption and internal thermal noise.

A critical hardware limitation of the ESP32 relevant to this project is
the relationship between the Wi-Fi radio and the second ADC group (ADC2).
When the Wi-Fi radio is active, ADC2 (which shares silicon with the RF
front-end hardware) becomes unavailable for analog measurements, returning
unreliable readings or causing crashes \citep{espressif2024}. This
mandates the use of ADC1 pins exclusively for analog current sensing, a
design constraint addressed in the hardware design of this project.

\section{IoT Communication Protocols}

\subsection{MQTT}

Message Queuing Telemetry Transport (MQTT) is a lightweight
publish-and-subscribe messaging protocol originally designed by IBM for
low-bandwidth, high-latency satellite communication links and subsequently
adopted by the OASIS standards body as an open protocol for IoT
applications \citep{oasis2019}. MQTT uses a central broker to route
messages between publishers and subscribers based on topic strings.
Publishers send messages without needing to know which subscribers exist;
subscribers receive all messages on topics they have registered interest
in. The protocol supports three Quality of Service levels governing
delivery guarantees, from fire-and-forget (QoS 0) through at-least-once
(QoS 1) to exactly-once (QoS 2).

MQTT's minimal header overhead (as little as two bytes for fixed header)
makes it significantly more efficient than HTTP for high-frequency sensor
telemetry. \citet{dzahir2023} confirmed that MQTT introduces substantially
lower communication overhead than HTTP in ESP32-based monitoring
applications, contributing to both lower energy consumption and reduced
latency. In this project, however, the primary cloud communication uses
Firebase HTTPS REST rather than a standalone MQTT broker, a design choice
discussed in the methodology.

\subsection{Bluetooth Low Energy}

Bluetooth Low Energy (BLE) is a wireless communication standard designed
for short-range, low-power data exchange. BLE uses the Generic Attribute
Profile (GATT) to organise data into services and characteristics,
allowing a client device to discover and subscribe to specific data streams
from a peripheral. The JK BMS broadcasts its telemetry as BLE GATT
notifications on a proprietary service UUID, transmitting 300-byte binary
frames fragmented across multiple 20-byte BLE notification packets
\citep{ijcrt2025}. The first ESP32 board acts as a BLE GATT client,
connecting to the BMS and subscribing to the notification characteristic
to receive these frames.

\section{Load Management in Energy Systems}

Intelligent load management involves the automated or semi-automated
control of electrical loads to optimise energy usage based on available
supply, demand priorities and system constraints. In the context of
battery-backed inverter systems, load management typically involves
disconnecting lower-priority loads before the battery reaches a critical
discharge level, thereby extending the runtime available to higher-priority
loads \citep{divya2009}.

Priority-based load shedding assigns each connected load a criticality
level. As the battery depletes, loads are shed in reverse order of
priority. Non-essential loads are disconnected first, preserving energy
for major loads. If depletion continues, major loads are shed next,
leaving only critical loads powered until the battery reaches the absolute
safe discharge limit. \citet{ijpeds2023} demonstrated that a
State-of-Charge-triggered load shedding system implemented on an ESP32
with relay switching achieved reliable and repeatable disconnection
sequences across extended discharge cycles.

\section{Firebase Realtime Database and Supabase in IoT Systems}

Firebase Realtime Database is a cloud-hosted, JSON-structured database
provided by Google that delivers data to connected clients over WebSocket
in real time \citep{expodocs2024}. Any write to the database is instantly
pushed to all clients that have registered listeners on the affected path,
without polling. This characteristic makes Firebase Realtime Database
particularly suitable for live telemetry dashboards where data changes
frequently and display latency must be minimised.

Supabase is an open-source Backend-as-a-Service platform built on
PostgreSQL, providing a REST API (PostgREST) for structured data storage
and querying \citep{firebasedocs2024}. Unlike Firebase's document model,
Supabase's relational schema and SQL querying capabilities make it better
suited for time-series analytics where complex aggregations, range queries
and multi-parameter filtering are required. In this project the two
platforms serve complementary roles: Firebase handles the live state
displayed on the dashboard, while Supabase stores the historical telemetry
used to generate analytics charts.

\section{Overview of Reviewed Literature}

Table \ref{tab:litoverview} presents a summary of the key literature reviewed in this chapter, highlighting the focus, methodology, and principal findings of each work and its relevance to the present project.

\begin{longtable}{L{2.8cm} L{3.5cm} L{4.5cm} L{3.5cm}}
\caption{Overview of Reviewed Literature} \label{tab:litoverview} \\
\toprule
\textbf{Author(s) \& Year} & \textbf{Topic / Focus} & \textbf{Key Findings} & \textbf{Relevance to This Project} \\
\midrule
\endfirsthead
\multicolumn{4}{c}{\tablename\ \thetable{} -- continued} \\
\toprule
\textbf{Author(s) \& Year} & \textbf{Topic / Focus} & \textbf{Key Findings} & \textbf{Relevance to This Project} \\
\midrule
\endhead
\midrule
\multicolumn{4}{r}{} \\
\endfoot
\bottomrule
\endlastfoot

Aliyu, Ramli \& Saleh (2013) & Nigeria electricity crisis and generation capacity & Documented poor grid reliability and low per-capita consumption despite installed capacity & Justifies the need for backup power systems in Nigeria \\

Ashton (2009) & Internet of Things concept & Introduced the foundational IoT paradigm of embedding connectivity in everyday objects & Provides the theoretical basis for connecting the inverter to the internet \\

Dini \& Ria (2024) & Li-ion battery characteristics and BMS integration & LiFePO\textsubscript{4} offers superior thermal stability and safety compared to other chemistries & Informs the battery chemistry selection for this project \\

Divya \& Ostergaard (2009) & Battery energy storage in power systems & Load management extends available runtime by disconnecting lower-priority loads first & Supports the priority-based load shedding approach \\

Dzahir \& Chia (2023) & ESP32 energy consumption in MQTT IoT systems & Deep sleep extends battery runtime significantly; MQTT has lower overhead than HTTP & Informs the dual-ESP32 architecture and cloud protocol choices \\

Espressif Systems (2024) & ESP32 technical reference manual & ADC2 is unavailable when Wi-Fi is active; ADC1 must be used for analog sensing during Wi-Fi operation & Directly informs the ADC1-only current sensor pin assignment \\

Expo (2024) & Using Firebase with Expo & Firebase Realtime Database provides instant data push to all connected clients & Supports the choice of Firebase for live dashboard state \\

Firebase (2024) & Cloud Firestore documentation & Document model suits real-time state synchronisation; relational backends suit analytics & Justifies the dual-backend architecture (Firebase + Supabase) \\

Fraunhofer ISE (2022) & Mechanical behaviour of prismatic cells & Unconstrained cycling causes electrode delamination; compression extends cycle life & Validates the use of steel compression plates in the battery pack \\

Hannan et al. (2017) & BMS technologies and energy management & Accurate SoC estimation is critical; coulomb counting is common but accumulates error & Informs the understanding of BMS SoC estimation methods \\

IJCRT (2025) & IoT-based battery monitoring with BMS and ESP32 & JK BMS transmits 300-byte binary frames fragmented over 20-byte BLE packets & Directly supports the BLE frame reassembly firmware design \\

IJPEDS (2023) & SoC-based intelligent load management for inverters & ESP32-based relay switching achieved reliable disconnection sequences across discharge cycles & Validates the feasibility of the proposed load shedding engine \\

KIT (2024) & Compression effects on LiFePO\textsubscript{4} cycle life & Appropriate compression extends cycle life; excessive compression causes degradation & Informs the compression fixture design for the battery pack \\

Nitta et al. (2015) & Li-ion battery materials present and future & LiFePO\textsubscript{4} uses iron-phosphate cathode and graphite anode; nominal voltage is 3.2V per cell & Provides the electrochemical foundation for the 4S pack design \\

OASIS (2019) & MQTT Version 5.0 specification & Publish-subscribe model with three QoS levels; minimal header overhead & Informs the communication protocol literature review \\

Omar et al. (2014) & LiFePO\textsubscript{4} aging and cycle life model & Repeated deep discharge permanently damages cells and reduces capacity & Justifies the need for protected load management to prevent deep discharge \\

Pillet et al. (2025) & Long-term cycling of LiFePO\textsubscript{4} cells & Capacity retention above 80\% after 2000 cycles under nominal conditions & Validates the longevity expectations for the chosen battery chemistry \\

Piller, Perrin \& Jossen (2001) & State-of-charge determination methods & Coulomb counting, OCV and model-based methods each have trade-offs in accuracy and complexity & Informs the understanding of SoC estimation limitations \\

Rashid (2014) & Power electronics handbook & Inverters convert DC battery energy to AC for appliances; battery-backed systems are silent and clean & Provides the power electronics foundation for the inverter subsystem \\

\end{longtable}

\section{Gaps in Existing Literature}

A review of the available literature reveals that while individual
components of the proposed system have been studied extensively, very few
works combine all of them into a single, hardware-verified implementation.
Most BMS interface studies use either Bluetooth-only or Wi-Fi-only
approaches, rarely addressing the fundamental incompatibility between
simultaneous BLE and Wi-Fi operation on a single ESP32 and the resulting
need for a distributed dual-microcontroller architecture. Works on load
prioritization typically simulate rather than implement the relay switching
hardware. IoT dashboard studies seldom implement the complete pipeline
from physical sensor reading through firmware processing to cloud
synchronisation and user-facing web application in a single project.

Table \ref{tab:litgaps} summarises the specific gaps identified in the
reviewed literature and indicates how this project addresses each of them.

{\renewcommand{\arraystretch}{1.5}
\begin{longtable}{L{4.0cm} L{5.5cm} L{5.3cm}}
\caption{Identified Gaps in Reviewed Literature and How This Project Addresses Them} \label{tab:litgaps} \\
\toprule
\textbf{Gap Identified} & \textbf{Limitation in Existing Work} & \textbf{How This Project Addresses the Gap} \\
\midrule
\endfirsthead
\multicolumn{3}{c}{\tablename\ \thetable{} (continued)} \\
\toprule
\textbf{Gap Identified} & \textbf{Limitation in Existing Work} & \textbf{How This Project Addresses the Gap} \\
\midrule
\endhead
\midrule
\multicolumn{3}{r}{} \\
\endfoot
\bottomrule
\endlastfoot

Simultaneous BLE and Wi-Fi operation & Most BMS interface studies use either Bluetooth-only or Wi-Fi-only approaches, rarely addressing the radio conflict on a single ESP32 & Implements a distributed dual-ESP32 architecture separating BLE and Wi-Fi responsibilities across two boards \\

Hardware-verified load prioritization & Works on load prioritization typically simulate rather than physically implement relay switching hardware & Implements and physically tests a relay-based load shedding engine with SoC thresholds and hysteresis \\

End-to-end IoT pipeline & IoT dashboard studies seldom implement the complete pipeline from sensor reading through firmware to cloud sync and web application & Delivers a fully working pipeline from BLE sensor acquisition to firmware processing, cloud publishing and a live Progressive Web Application \\

JK BMS protocol documentation & Few accessible references document the JK BMS proprietary binary frame structure for firmware-level integration & Provides a working firmware-level parser for the JK02\_32S binary protocol as a practical reference \\

ADC/Wi-Fi hardware conflict handling & ESP32-based IoT designs commonly overlook the ADC2/Wi-Fi radio conflict during analog sensing & Explicitly routes all current sensing to ADC1-only pins to avoid conflict during active Wi-Fi operation \\

Dual-backend cloud architecture & Existing systems typically rely on a single cloud backend, limiting either live-state responsiveness or historical analytics capability & Combines Firebase Realtime Database for live state with Supabase for structured historical analytics \\

\end{longtable}

This project addresses all of these gaps simultaneously by implementing a
complete, working system verified on physical hardware. The dual-ESP32
architecture, the JK BMS binary protocol parser, the ADC1-only constraint
for current sensing during Wi-Fi operation, the dual-backend cloud
architecture combining Firebase and Supabase, and the Progressive Web
Application interface together represent a contribution to the practical
literature on IoT-enabled energy management systems.



%%======================================================
%%  CHAPTER 3: METHODOLOGY
%%======================================================
\newpage
\chapter{METHODOLOGY}

\section{System Overview}

The complete system is an embedded firmware ecosystem for real-time
monitoring, intelligent load management and cloud telemetry for a
lithium-ion battery pack managed by a JK Battery Management System. The
system architecture is illustrated in Figure~\ref{fig:sysarch}.

\begin{figure}[H]
\centering
\begin{tikzpicture}[
  scale=1,
  transform shape,
  node distance=1.0cm
]
  \node[sysblock, fill=blockgreen] (batt)
    {LiFePO$_4$ Battery Pack};

  \node[sysblock, fill=blockgreen, below=of batt] (bms)
    {JK-BMS};

  \node[sysblock, fill=blockblue, below=of bms] (esp1)
    {ESP32 Board 1 (BLE + TFT)};

  \node[sysblock, fill=blockblue, below=of esp1] (esp2)
    {ESP32 Board 2 (Wi-Fi + Relay)};

  \node[sysblock, fill=blockgrey, below=of esp2] (router)
    {Wi-Fi Router / 4G MiFi};

  \node[sysblock, fill=blockorange, below=of router] (cloud)
    {Firebase RTDB + Supabase};

  \node[sysblock, fill=blockblue, below=of cloud] (app)
    {Progressive Web App};

  \draw[arrow] (batt) -- node[right,font=\tiny]{DC} (bms);
  \draw[dasharrow] (bms) -- node[right,font=\tiny]{BLE} (esp1);
  \draw[arrow] (esp1) -- node[right,font=\tiny]{UART} (esp2);
  \draw[arrow] (esp2) -- node[right,font=\tiny]{Wi-Fi HTTPS} (router);
  \draw[arrow] (router) -- node[right,font=\tiny]{Internet} (cloud);
  \draw[arrow] (cloud) -- node[right,font=\tiny]{WebSocket / REST} (app);

  \draw[dasharrow] (app.east) to[bend right=45]
      node[right,font=\tiny]{Commands} (cloud.east);

\end{tikzpicture}
\caption{Overall System Architecture}
\label{fig:sysarch}
\end{figure}



The reason for splitting responsibilities across two microcontrollers is
a real hardware limitation: Bluetooth Low Energy and Wi-Fi cannot run
reliably at the same time on a single ESP32. When both radios are active
on one board, the result is severe packet loss on the BLE link, BMS
connection timeouts and stuttering on the TFT display. Separating the
roles solves this cleanly. ESP32 Board 1 owns the BLE connection to the
BMS and the TFT display. ESP32 Board 2 owns everything network-related:
Wi-Fi, relay switching, current sensing, Firebase and Supabase.

\section{Hardware Design}

The hardware design described in this section directly addresses
Objective~(i), which calls for the design and assembly of a 4-series
LiFePO\textsubscript{4} battery pack and the integration of the JK Battery
Management System for cell-level monitoring and protection.

\subsection{Battery Pack Design and Assembly}

The battery pack consists of four 100\,Ah 3.2\,V nominal Lithium Iron
Phosphate ($\text{LiFePO}_4$) prismatic cells connected in series, producing a 4S configuration
with a nominal pack voltage of 12.8\,V and a full-charge voltage of 14.6\,V.
The series connection increases the total voltage while maintaining the
100\,Ah capacity of a single cell, giving the pack approximately 1.28\,kWh
of usable energy storage. Figure~\ref{fig:cells_purchased} shows the four individual 100\,Ah prismatic cells procured for the system before physical integration.

\begin{figure}[H]
  \centering
  \includegraphics[
    width=0.55\textwidth,
    angle=90
  ]{image of the lifepo4 battery purchased_.jpg}
  \caption{The four 100\,Ah LiFePO\textsubscript{4} prismatic cells purchased for the project}
  \label{fig:cells_purchased}
\end{figure}

Before assembly, each cell was individually tested with a digital multimeter
to verify voltage uniformity. Cells with significant voltage deviations,
physical swelling or visible damage were excluded. The cells were then
top-balanced by charging each one individually to 3.65\,V before connection,
ensuring all four entered service at the same State of Charge. This gives
the BMS coulomb counter an accurate starting reference.

The four cells were arranged in alternating orientations so that adjacent
positive and negative terminals faced each other, enabling a compact series
connection using copper busbars. Steel end plates were positioned on both
outer faces of the cell stack, and four threaded steel rods were passed
through alignment holes at the plate corners. Bolts and nuts were tightened
to apply controlled mechanical compression across the cell faces.

This compression is essential for preventing electrode delamination. During
each charge and discharge cycle, the intercalation and de-intercalation of
lithium ions causes the electrode layers to expand and contract. Without
constraint, this repeated volumetric change causes progressive separation of
electrode layers from each other and from their current collectors, which
raises internal resistance and reduces capacity over time. The steel plate
fixture maintains intimate electrode contact throughout the lifetime of the
pack \citep{fraunhofer2022}.

\subsection{Physical Pack Construction: Enclosure, Insulation and Compression}

\subsubsection{Inter-Cell Electrical Insulation}

To prevent accidental electrical short circuits across adjacent metallic cell bodies, custom 3D-printed insulating spacers were engineered and placed between every adjacent cell pair and between the cell stack and the metallic enclosure walls. Because the prismatic cells possess conductive aluminium outer walls, direct physical contact across adjacent cell bodies or with the metallic chassis would establish inadvertent conductive paths. Figure~\ref{fig:spacers} illustrates the 3D-printed polymer spacers that eliminate this risk by providing a continuous mechanical and electrical barrier at every contact boundary.

\begin{figure}[H]
  \centering
  \includegraphics[width=0.60\textwidth]{image of the 3d printed spacers that was but in between the battery cells to prevent short-circuit between battery cells.png}
  \caption{3D-printed polymer spacers placed between the LiFePO\textsubscript{4} cells to prevent inter-cell short circuits}
  \label{fig:spacers}
\end{figure}

\subsubsection{Compression Plates}

Prismatic LiFePO\textsubscript{4} cells experience cyclic volumetric expansion (``breathing'') during intercalation and de-intercalation of lithium ions. Without rigid mechanical restraint, this cyclic swelling leads to internal delamination of the electrode and separator layers, casing distortion, and premature capacity loss. To counteract this degradation mechanism, steel compression end-plates and 8\,mm threaded steel rods were fabricated to clamp the four cells under continuous, uniform mechanical pressure. Figure~\ref{fig:compression_parts} shows the engineering drawing and hardware components of the compression fixture, while Figure~\ref{fig:assembly_sequence} illustrates the pack assembly sequence.

\begin{figure}[H]
  \centering
  \begin{subfigure}[b]{0.48\textwidth}
    \centering
    \includegraphics[width=\textwidth]{engineering drawing of the battery compression plate.png}
    \caption{Engineering drawing of the steel compression plate}
    \label{fig:plate_drawing}
  \end{subfigure}
  \hfill
  \begin{subfigure}[b]{0.48\textwidth}
    \centering
    \includegraphics[width=\textwidth]{image of the compression plates, the 8mm threaded rods and the 4 nuts used for the battery pack compression.png}
    \caption{Steel compression plates, 8\,mm threaded rods and nuts}
    \label{fig:compression_hardware}
  \end{subfigure}
  \caption{Battery pack compression fixture components}
  \label{fig:compression_parts}
\end{figure}

\begin{figure}[htbp]
  \centering
  \begin{subfigure}[b]{0.48\textwidth}
    \centering
    \includegraphics[width=\textwidth]{image of the 4 lifepo4 cells, with the 3d printed spacers in between them and the compression plates at either end.jpg}
    \caption{Cells arranged with spacers and compression plates before tightening}
    \label{fig:cells_assembled}
  \end{subfigure}
  \hfill
  \begin{subfigure}[b]{0.48\textwidth}
    \centering
    \includegraphics[width=\textwidth]{assembling and compression of the battery and BMS integration_.jpg}
    \caption{Battery pack assembly and BMS integration in progress}
    \label{fig:assembly_progress}
  \end{subfigure}
  \caption{Battery pack assembly sequence}
  \label{fig:assembly_sequence}
\end{figure}

\subsection{Electrical Interconnection: Cabling, Lugs and Torque Practice}
\label{sec:cablesizing}

The main DC power connections between the battery pack terminals, the BMS switching path, and the inverter DC input were executed using high-flexibility $16\,\text{mm}^{2}$ multi-strand copper cable. The selection of $16\,\text{mm}^{2}$ conductor cross-section was determined through rigorous electrical engineering calculations satisfying continuous ampacity, thermal dissipation, and permissible voltage drop criteria.

\paragraph{1. Maximum Design Current Calculation.}
The inverter employed is rated at an apparent power $S = 1.5\,\text{kVA} = 1500\,\text{VA}$. Assuming a standard inductive power factor $\cos\phi = 0.8$, the maximum real output power is:
\begin{equation}
    P_{\text{ac}} = S \times \cos\phi = 1500\,\text{VA} \times 0.8 = 1200\,\text{W}
    \label{eq:pac}
\end{equation}
Accounting for the inverter's DC-to-AC conversion efficiency ($\eta \approx 88\% = 0.88$), the continuous DC power demand drawn from the battery pack is:
\begin{equation}
    P_{\text{dc}} = \frac{P_{\text{ac}}}{\eta} = \frac{1200\,\text{W}}{0.88} \approx 1363.6\,\text{W}
    \label{eq:pdc}
\end{equation}
The maximum continuous DC discharge current occurs at the minimum operating pack voltage before undervoltage cutoff. Setting a worst-case threshold of $V_{\text{min}} = 11.2\,\text{V}$ (equivalent to $2.80\,\text{V}$ per cell, below which the BMS cuts off at $3.00\,\text{V}$):
\begin{equation}
    I_{\text{dc,max}} = \frac{P_{\text{dc}}}{V_{\text{min}}} = \frac{1363.6\,\text{W}}{11.2\,\text{V}} \approx 121.75\,\text{A}
    \label{eq:idc_max}
\end{equation}
Under nominal operating voltage ($V_{\text{nom}} = 12.8\,\text{V}$) and continuous practical test loads ($800\,\text{W}$ to $1000\,\text{W}$), the continuous discharge current is:
\begin{equation}
    I_{\text{dc,nom}} = \frac{1000\,\text{W} / 0.88}{12.8\,\text{V}} \approx 88.8\,\text{A}
    \label{eq:idc_nom}
\end{equation}
During charging, the maximum configured inverter charge current is limited to $I_{\text{chg}} = 30.0\,\text{A}$.

\paragraph{2. Current-Carrying Capacity (Ampacity) and Thermal Rating.}
According to IEC 60364-5-52 and BS 7671 standards for multi-strand flexible copper conductors with high-temperature silicone insulation in free air and enclosed trunking, a $16\,\text{mm}^{2}$ copper cable has a baseline continuous ampacity of $I_z \approx 110\,\text{A}$ at an ambient temperature of $30\,^{\circ}\text{C}$. Applying a thermal derating correction factor $k_t = 0.94$ for tropical ambient conditions up to $40\,^{\circ}\text{C}$ in Akure:
\begin{equation}
    I_z' = I_z \times k_t = 110\,\text{A} \times 0.94 \approx 103.4\,\text{A}
    \label{eq:ampacity}
\end{equation}
This rating comfortably exceeds the continuous operational current demands of the system under all sustained load and charging scenarios.

\paragraph{3. Conductor Loop Resistance and Voltage Drop.}
For low-voltage ($12\,\text{V}$) DC power systems, minimizing percentage voltage drop is critical to avoid premature undervoltage tripping and power loss. The total conductor loop length (positive and negative leads combined) between the battery pack, BMS, and inverter terminals is $L_{\text{loop}} = 1.0\,\text{m}$ ($0.5\,\text{m}$ per lead).

The resistivity of annealed copper at its maximum conductor operating temperature ($75\,^{\circ}\text{C}$) is $\rho_{75} = 2.14 \times 10^{-8}\,\Omega\cdot\text{m}$. For a cross-sectional area $A = 16\,\text{mm}^{2} = 16 \times 10^{-6}\,\text{m}^{2}$, the total cable loop resistance is:
\begin{equation}
    R_{\text{cable}} = \frac{\rho_{75} \cdot L_{\text{loop}}}{A} = \frac{2.14 \times 10^{-8}\,\Omega\cdot\text{m} \times 1.0\,\text{m}}{16 \times 10^{-6}\,\text{m}^{2}} = 1.3375 \times 10^{-3}\,\Omega = 1.338\,\text{m}\Omega
    \label{eq:rcable}
\end{equation}
At a high continuous operational discharge current of $I = 80\,\text{A}$, the total voltage drop across the DC cabling is:
\begin{equation}
    \Delta V = I \cdot R_{\text{cable}} = 80\,\text{A} \times 1.3375 \times 10^{-3}\,\Omega \approx 0.107\,\text{V}
    \label{eq:vdrop}
\end{equation}
Expressed as a percentage of the nominal $12.8\,\text{V}$ battery bus:
\begin{equation}
    \%\Delta V = \left(\frac{\Delta V}{V_{\text{nom}}}\right) \times 100\% = \left(\frac{0.107\,\text{V}}{12.8\,\text{V}}\right) \times 100\% \approx 0.836\%
    \label{eq:pvdrop}
\end{equation}
This voltage drop of $0.836\%$ is significantly below the standard $3.0\%$ maximum limit prescribed by IEEE and IEE Wiring Regulations for DC power circuits.

\paragraph{4. Ohmic Power Loss.}
The total thermal power dissipated along the cable length at $80\,\text{A}$ is:
\begin{equation}
    P_{\text{loss}} = I^{2} \cdot R_{\text{cable}} = (80\,\text{A})^{2} \times 1.3375 \times 10^{-3}\,\Omega \approx 8.56\,\text{W}
    \label{eq:ploss}
\end{equation}
This modest dissipation ($<8.6\,\text{W}$ distributed over a 1\,m cable span) generates negligible temperature rise, verifying that $16\,\text{mm}^{2}$ copper cable delivers superior safety, high electrical efficiency, and minimal thermal stress.

Cable lugs were hydraulically crimped onto every cable termination point: the cell terminal busbars, the BMS main power terminals, and the inverter DC input studs. Bare stripped conductors clamped under screw heads are unsuitable for high currents due to inconsistent contact area and vulnerability to creep. A crimped copper lug provides a large, uniform interfacial contact area. Connections were tightened with a calibrated torque wrench to ensure tight, low-resistance mechanical joints ($<0.05\,\text{m}\Omega$ contact resistance), eliminating hotspot formation.

\subsection{JK Battery Management System}

The protection, cell balancing, and telemetry acquisition for the 4S pack are overseen by the smart JK Battery Management System (BMS). The BMS unit is illustrated in Figure~\ref{fig:jkbms}, while Figure~\ref{fig:bms_integration} shows its physical wiring and final placement inside the metallic battery enclosure.

\begin{figure}[H]
  \centering
  \includegraphics[width=0.70\textwidth, trim=0cm 6cm 2cm 1cm,
       clip]{image of the jk bms.jpg}
  \caption{The JK Battery Management System unit used in this project}
  \label{fig:jkbms}
\end{figure}

\begin{figure}[H]
  \centering

  \begin{subfigure}[b]{0.4\textwidth}
    \centering
    \includegraphics[
      width=\textwidth,
      angle=90,
      height=9cm
    ]{image of the assembled battery with BMS wire connections on it.jpg}
    \caption{BMS balance leads connected to the battery pack terminals}
    \label{fig:bms_wired}
  \end{subfigure}
  \hfill
  \begin{subfigure}[b]{0.4\textwidth}
    \centering
    \includegraphics[
      width=\textwidth,
       height=9cm,
       trim=3cm 6cm 2cm 1cm,
       clip
    ]{image of the assembled battery in the metal casing container_.jpg}
    \caption{Fully assembled battery pack housed in the metallic enclosure}
    \label{fig:pack_cased}
  \end{subfigure}

  \caption{BMS integration and final battery pack assembly}
  \label{fig:bms_integration}
\end{figure}

The BMS is connected to the battery pack in two ways. First, a set of thin
cell-tap sense wires, one per inter-cell junction plus the pack positive and
negative terminals, allows the BMS to measure each cell's voltage
individually. Second, the main 16\,mm$^{2}$ power leads carry the full pack
charge and discharge current through the BMS's internal MOSFET switching
path. This series placement is what allows the BMS to physically disconnect
the battery from the load or charger when a protection threshold is exceeded.

The BMS protection and charge-control parameters were configured via the
manufacturer's Bluetooth mobile application. Table~\ref{tab:bmsconfig} shows
the configured values.

\begin{table}[htbp]
\centering
\caption{JK BMS Protection Configuration Parameters}
\label{tab:bmsconfig}
\begin{tabular}{L{5.5cm} C{3.0cm} C{3.5cm}}
\toprule
\textbf{Parameter} & \textbf{Configured Value} & \textbf{Pack Equivalent (4S)} \\
\midrule
Cell over-voltage protection (COVP) & 3.55\,V & 14.20\,V \\
Cell under-voltage protection (CUVP) & 3.00\,V & 12.00\,V \\
Balance trigger voltage & 3.36\,V & 13.44\,V \\
Float charge voltage & 3.375\,V avg & 13.50\,V \\
\bottomrule
\end{tabular}
\end{table}

The configured parameters, shown in the BMS application interface in Figure~\ref{fig:bms_settings}, maintain a deliberate safety margin within the chemistry's absolute electro-chemical boundaries of $2.50\,\text{V}$ to $3.65\,\text{V}$ per cell. This conservative strategy sacrifices a negligible fraction ($<2\%$) of extreme-end capacity in exchange for drastically reduced cell degradation and extended cyclic lifespan.

\begin{figure}[H]
  \centering
  \includegraphics[width=0.45\textwidth]{bms_settings.png}
  \caption{BMS mobile application configuration page showing protection and charge parameters}
  \label{fig:bms_settings}
\end{figure}

\subsection{Inverter Integration and Charge-Profile Configuration}

The energy conversion and primary battery charging are performed by a Nexus 1.5\,kVA hybrid solar inverter. Its rear interface panel, depicted in Figure~\ref{fig:inverter_rear}, features AC input/output terminals (L and N), a high-efficiency PV MPPT input terminal pair (rated $18\,\text{V}$--$30\,\text{V}$ optimum), heavy-duty DC $12\,\text{V}$ battery terminal lugs, an inverter power switch, and a thermostatically controlled cooling fan. Table~\ref{tab:inverterspec} details the inverter electrical specifications and the charge parameters configured for the system.

\begin{figure}[H]
  \centering
  \includegraphics[width=0.65\textwidth]{the rear of the inverter showing the ac input and output, pv input, DC 12V input, a switch to on and off the inverter, etc.jpg}
  \caption{Rear panel of the Nexus 1.5\,kVA hybrid inverter showing AC terminals, PV input, DC battery input and cooling fan}
  \label{fig:inverter_rear}
\end{figure}

\begin{table}[H]
\centering
\caption{Nexus 1.5\,kVA Hybrid Inverter Specifications and Configuration}
\label{tab:inverterspec}
\begin{tabular}{L{5.5cm} L{4.5cm} L{3.0cm}}
\toprule
\textbf{Parameter} & \textbf{Value} & \textbf{Source} \\
\midrule
Rated capacity & 1.5\,kVA & Project specification \\
Battery input & DC\,12V (nominal) & Rear panel label \\
PV input range & 16--80\,V (best 18--30\,V) & Rear panel label \\
AC input/output & L/N terminals, input and output & Rear panel label \\
Bulk charge voltage (configured) & 14.5\,V & Inverter settings \\
Charge current limit (configured) & 30\,A & Inverter settings \\
\bottomrule
\end{tabular}
\end{table}

One coordination point worth noting: the inverter's configured bulk charge
target of 14.5\,V corresponds to an average of 3.625\,V per cell. This is
higher than the BMS's configured COVP of 3.55\,V per cell, which is equivalent
to 14.20\,V at the pack level. In practice the BMS cell-level protection is
the more conservative limit and is expected to begin restricting charge
current before the inverter's 14.5\,V bulk target is reached uniformly. This
is a desirable coordination outcome: the inverter's charging profile is not
the final arbiter of pack safety, and the BMS's own per-cell cutoff is the
effective limiting factor.

\subsection{BMS Calibration Procedure and State of Charge Dynamics}
\label{sec:bmscalib}

Following mechanical and electrical assembly, the battery pack and BMS were connected to the inverter to execute an initial full-charge calibration routine under supervisory guidance. In Lithium Iron Phosphate ($\text{LiFePO}_4$) chemistry, the open-circuit voltage curve is extraordinarily flat across $10\%$ to $90\%$ SoC, exhibiting a voltage variation of only a few tens of millivolts (around $3.25\,\text{V}$ to $3.35\,\text{V}$ per cell). Consequently, voltage-based SoC estimation is fundamentally unreliable in this plateau region, requiring the BMS to rely on coulomb counting (current integration over time). However, uncalibrated coulomb counters accumulate drift and have no initial absolute reference point.

During the initial commissioning charge, the BMS attempts to locate its true $100\%$ SoC datum point. Figure~\ref{fig:bms_charging} captures a critical diagnostic observation during this initial calibration charge.

\begin{figure}[H]
  \centering
  \includegraphics[width=0.3\textwidth]{the jk bms interface showing the battery percentage, charging current, voltage diff, max cell voltage, power, temperature, etc.jpg}
  \caption{JK BMS mobile application showing live battery metrics during the initial calibration charge}
  \label{fig:bms_charging}
\end{figure}

In Figure~\ref{fig:bms_charging}, the BMS application circular gauge indicates an apparent State of Charge of $99\%$. However, examination of the live electrical metrics reveals that the pack is actively absorbing a substantial charging current of $26.75\,\text{A}$ ($359.8\,\text{W}$) at a pack terminal voltage of only $13.45\,\text{V}$ (an average of $3.362\,\text{V}$ per cell). In a genuinely full $\text{LiFePO}_4$ battery ($99\%$--$100\%$ SoC), the cell operating points enter the steep upper voltage ``knee'' ($>3.50\,\text{V}$ per cell, or $>14.0\,\text{V}$ pack voltage), causing the constant-voltage absorption current to rapidly taper down toward zero. The fact that the battery in Figure~\ref{fig:bms_charging} is still accepting nearly $27\,\text{A}$ of bulk charge current at $13.45\,\text{V}$ proves conclusively that the pack is still in the middle of its bulk absorption plateau and is physically nowhere near $99\%$ true capacity.

This discrepancy highlights the necessity of the calibration procedure: charging was sustained at constant voltage until the pack voltage reached $14.20\,\text{V}$ and the charging current naturally tapered down to a tail current of approximately $2.0\,\text{A}$ ($\approx\text{C}/50$). Only at this low tail current does the BMS recalibrate and lock its internal $100\%$ State of Charge reference, eliminating coulomb counter drift for all subsequent discharge cycles.

\subsection{Active Cell Balancing and Differential Voltage Verification}

Active cell balancing is a vital feature of the JK BMS. Unlike conventional passive balancing systems that dissipate excess energy from higher-voltage cells as wasted heat through bleed resistors, the JK BMS incorporates an active bidirectional inductive balancing circuit \citep{hannan2017}. This circuit dynamically shuttles charge from the highest-voltage cell directly into the lowest-voltage cell at balancing currents up to $1.0\,\text{A}$ (measured at $0.970\,\text{A}$ during active equalization).

Figure~\ref{fig:bms_status} presents the detailed cell diagnostic status captured when the pack is resting ($12.77\,\text{V}$ pack voltage, $0.00\,\text{A}$ current, $61\%$ SoC).

\begin{figure}[H]
  \centering
  \includegraphics[width=0.4\textwidth, scale=0.7]{bmsactivebal.png}
  \caption{JK BMS status screen showing individual cell voltages, cell wire resistances, temperatures and differential voltage}
  \label{fig:bms_status}
\end{figure}

The individual cell voltages in Figure~\ref{fig:bms_status} are recorded as $\text{Cell}_1 = 3.193\,\text{V}$, $\text{Cell}_2 = 3.191\,\text{V}$, $\text{Cell}_3 = 3.190\,\text{V}$, and $\text{Cell}_4 = 3.193\,\text{V}$. The critical diagnostic metric is the Cell Voltage Difference ($\text{DV} = V_{\max} - V_{\min}$), which is measured at an exceptionally low value of:
\begin{equation}
    \text{DV} = 3.193\,\text{V} - 3.190\,\text{V} = 0.003\,\text{V} = 3.0\,\text{mV}
    \label{eq:celldv}
\end{equation}
A differential voltage of only $3.0\,\text{mV}$ across all four series cells demonstrates that the four prismatic cells are closely matched in electrochemical capacity, state of charge, and internal resistance. Furthermore, the measured cell wire resistances are uniform across all channels ($0.051\,\Omega$ to $0.053\,\Omega$), confirming that the balance leads and terminal lugs exhibit consistent contact resistance. High cell consistency ensures that all cells reach charge and discharge boundaries simultaneously, preventing any single weak cell from prematurely triggering protective cutoffs and maximizing the pack's usable energy extraction.

\subsection{Key Battery and BMS Terminology}

The following terms recur throughout this report and are defined here for
clarity.

\begin{enumerate}[label=\Alph*., leftmargin=0.7in, noitemsep]
  \item \textbf{State of Charge (SoC):} the battery's remaining charge
        expressed as a percentage of rated capacity.
  \item \textbf{Depth of Discharge (DoD):} the complement of SoC; the
        fraction of rated capacity that has been discharged from full.
  \item \textbf{State of Health (SoH):} remaining capacity or performance
        relative to the original as-new specification.
  \item \textbf{C-rate:} current expressed relative to rated capacity. 1C
        equals 100\,A for a 100\,Ah pack; C/50 equals 2\,A.
  \item \textbf{Nominal voltage:} the characteristic rated voltage under
        typical conditions (3.2\,V/cell; 12.8\,V for a 4S pack).
  \item \textbf{Cutoff voltage:} the voltage at which charging or
        discharging is halted to protect the cell.
  \item \textbf{Float voltage:} a maintenance charging voltage held after
        full charge to offset self-discharge without overcharging.
  \item \textbf{Bulk/Absorption/Float charging stages:} the three phases
        of a multi-stage charge profile. Bulk is constant current with
        rising voltage; absorption is constant voltage with tapering current;
        float is a lower maintenance voltage once the battery is full.
  \item \textbf{Tail current:} the low current at which absorption-phase
        charging is considered complete.
  \item \textbf{Cell balancing:} equalising charge across series-connected
        cells. Passive balancing wastes excess charge as heat; active
        balancing transfers it between cells.
  \item \textbf{Coulomb counting:} SoC estimation by integrating measured
        current over time relative to rated capacity.
  \item \textbf{Open-circuit voltage (OCV):} terminal voltage with no load
        or charge current flowing; used as a secondary SoC cross-check.
  \item \textbf{4S configuration:} four cells connected in series; voltages
        add while capacity stays equal to one cell.
  \item \textbf{Busbar:} a rigid conductive strip making a low-resistance
        connection between cell terminals.
  \item \textbf{Cable lug:} a crimped terminal providing a bolt-compatible
        connection at high-current joints.
  \item \textbf{Ampacity:} the maximum continuous current a conductor can
        carry without exceeding its temperature rating.
  \item \textbf{Compression plate:} a rigid end-plate clamped across a cell
        stack to apply constant pressure and resist swelling.
  \item \textbf{Thermal runaway:} an uncontrolled self-accelerating
        temperature rise within a cell triggered by severe overcharge,
        internal short circuit or physical damage.
\end{enumerate}

\subsection{ESP32 Board 1: BMS Gateway and TFT Display}

ESP32 Board 1 serves two functions: it maintains the BLE connection to the
JK BMS and drives the 2.4-inch ILI9341 TFT display over SPI.
Table~\ref{tab:esp1pins} shows the complete pin assignment.

\begin{longtable}{L{3.2cm} L{3.8cm} C{2.0cm} L{4.0cm}}
\caption{ESP32 Board 1 Pin Assignment}
\label{tab:esp1pins}\\

\toprule
\textbf{Terminal} & \textbf{Connection} & \textbf{GPIO} & \textbf{Description} \\
\midrule
\endfirsthead

\multicolumn{4}{c}%
{{\tablename\ \thetable{} -- continued }} \\
\toprule
\textbf{Terminal} & \textbf{Connection} & \textbf{GPIO} & \textbf{Description} \\
\midrule
\endhead

\midrule
\multicolumn{4}{r}{{}} \\
\endfoot

\bottomrule
\endlastfoot

TFT VCC & TFT VCC rail & 5V pin & 5V supply to TFT module \\
TFT LED & TFT backlight & 3.3V pin & 3.3V to TFT LED pin \\
TFT GND & TFT GND rail & GND pin & Common ground \\
TFT MISO & TFT SDO & 19 & SPI data from display \\
TFT MOSI & TFT SDA/SDI & 23 & SPI data to display \\
TFT SCLK & TFT SCL/CLK & 18 & SPI clock \\
TFT CS & TFT CS & 5 & SPI Chip Select (active LOW) \\
TFT DC & TFT DC/RS & 2 & Data/Command selector \\
TFT RST & TFT RST & 4 & Hardware reset \\
UART2 TX & Board 2 GPIO 16 & 17 & Sends BMS telemetry JSON \\
UART2 RX & Board 2 GPIO 17 & 16 & Receives heartbeat JSON \\
GND (UART) & Board 2 GND & GND pin & Shared UART ground \\

\end{longtable}



\subsection{ESP32 Board 2: Network, Relays and Current Sensors}

ESP32 Board 2 manages Wi-Fi connectivity, relay switching and analog current
measurement. Three ACS758 Hall-effect current sensors are connected to this
board, alongside three relay modules that control the three load output
channels. Table~\ref{tab:esp2pins} shows the pin assignment.

\begin{longtable}{L{3.2cm} L{3.8cm} C{2.0cm} L{4.0cm}}
\caption{ESP32 Board 2 Pin Assignment}
\label{tab:esp2pins}\\

\toprule
\textbf{Terminal} & \textbf{Connection} & \textbf{GPIO} & \textbf{Description} \\
\midrule
\endfirsthead

\multicolumn{4}{c}%
{{\tablename\ \thetable{} -- continued }} \\
\toprule
\textbf{Terminal} & \textbf{Connection} & \textbf{GPIO} & \textbf{Description} \\
\midrule
\endhead

\midrule
\multicolumn{4}{r}{{}} \\
\endfoot

\bottomrule
\endlastfoot

VCC (Sensors) & ACS758 VCC rail & 3.3V pin & Powers all three sensors \\
GND & Sensor and relay GND & GND pin & Common ground \\
UART2 RX & Board 1 GPIO 17 & 16 & Receives BMS telemetry JSON \\
UART2 TX & Board 1 GPIO 16 & 17 & Sends heartbeat JSON \\
RELAY CH1 & Relay module IN1 & 25 & Critical load relay (Active HIGH) \\
RELAY CH2 & Relay module IN2 & 26 & Major load relay (Active HIGH) \\
RELAY CH3 & Relay module IN3 & 27 & Non-essential load relay (Active HIGH) \\
SENSOR CH1 & ACS758 \#1 VIOUT & 32 & CH1 current (ADC1) \\
SENSOR CH2 & ACS758 \#2 VIOUT & 33 & CH2 current (ADC1) \\
SENSOR CH3 & ACS758 \#3 VIOUT & 34 & CH3 current (ADC1) \\
VCC (Relay) & Relay module VCC & 5V pin & Powers relay coil drivers \\
GND (Relay) & Relay module GND & GND pin & Relay ground \\

\end{longtable}

A critical hardware decision concerns which ADC pins the current sensors use.
The ESP32 has two ADC groups. ADC2 shares silicon with the Wi-Fi radio: when
Wi-Fi is active, ADC2 is completely unavailable and returns garbage readings.
GPIOs 32, 33 and 34 belong exclusively to ADC1 and remain fully functional
at all times regardless of Wi-Fi state. All three ACS758 outputs are therefore
wired to ADC1 pins only.

\begin{figure}[H]
  \centering
  \begin{subfigure}[b]{0.48\textwidth }
    \centering
    \includegraphics[width=\textwidth, angle=90, height=10cm]{The 4 channel relay 1.jpg}
    \caption{The four-channel relay module used for load switching}
    \label{fig:relay_module}
  \end{subfigure}
  \hfill
  \begin{subfigure}[b]{0.48\textwidth}
    \centering
    \includegraphics[width=\textwidth]{The three current sensors attached to the relay_.jpg}
    \caption{ACS758 current sensors attached to the relay output terminals}
    \label{fig:sensors_relay}
  \end{subfigure}
  \caption{Relay module and ACS758 current sensor hardware}
  \label{fig:relay_sensors}
\end{figure}

The relay modules use active-high logic. Sending a HIGH signal to a relay
input pin energises the relay coil and closes the Normally Open contact,
connecting the load. Sending LOW de-energises the coil and opens the circuit.
The relay COM terminals are looped together and connected to the inverter AC
output. Each Normally Open terminal connects to the IP+ terminal of its
corresponding ACS758 current sensor, so that load current flows through the
sensor whenever the relay is closed. This arrangement allows the current drawn
by each individual load channel to be measured independently.

\subsection{ACS758 Current Sensor Configuration}

The ACS758 is a Hall-effect current sensor. It measures current flowing
through a conductive path embedded in the sensor chip without any galvanic
contact between the high-current circuit and the low-voltage sensing
electronics. At zero current, the sensor output sits at approximately half
the supply voltage (VCC/2, roughly 1.65\,V on a 3.3\,V supply). The output
shifts upward for positive current and downward for negative current.

Auto-calibration is performed at firmware startup by reading the sensor output
with no load connected and storing that average reading as the zero-current
reference. This corrects for manufacturing offset variation from unit to unit.
Multi-sample averaging over 64 ADC readings reduces the impact of ADC noise,
and a deadband filter discards any reading below approximately 0.1\,A, which
prevents ADC noise from appearing as a phantom current when the relay is open.

\subsection{DC-DC Buck Converter for Controller Power}

The controller electronics (ESP32 microcontrollers, relay driver modules, ACS758 current sensors, and TFT display) are powered directly from the $12\,\text{V}$ battery bank via an efficient DC-DC step-down (buck) converter module, illustrated in Figure~\ref{fig:buck}. Because the battery terminal voltage fluctuates between $10.0\,\text{V}$ and $14.6\,\text{V}$ during dynamic charge and discharge cycles, the buck converter steps this variable DC input down to a strictly regulated $5.0\,\text{V}$ rail. The high efficiency ($>90\%$) of the buck topology minimizes parasitic quiescent power draw from the battery bank and avoids the significant thermal dissipation inherent to linear regulators.

\begin{figure}[H]
  \centering
  \includegraphics[width=0.45\textwidth, angle=90, trim= 14cm 0cm 0cm 0cm, clip]{The buck converter.jpg}
  \caption{DC-DC step-down (buck) converter supplying regulated 5\,V to the controller electronics from the 12\,V battery bus}
  \label{fig:buck}
\end{figure}

\section{Software Design}

The software design described in this section fulfills Objective~(ii), which
calls for firmware development for both ESP32 microcontrollers, and lays the
groundwork for Objectives~(iii) and~(iv) through the load shedding engine
and the cloud-connected web application.

\subsection{ESP32 Board 1 Firmware}

\subsubsection{BLE Communication with the JK BMS}

The firmware uses the NimBLE-Arduino library, a memory-optimised Bluetooth
stack for ESP32, to establish and maintain the BLE connection to the JK BMS.
NimBLE was chosen over the default ArduinoBLE stack because it uses
significantly less RAM, leaving more headroom for the JSON serialiser and
the TFT rendering buffer to operate concurrently.

The connection procedure is:
\begin{enumerate}[noitemsep, leftmargin=0.7in]
  \item The ESP32 initialises \texttt{BLEDevice} and creates a \texttt{BLEClient}.
  \item It connects directly to the BMS using its hardcoded MAC address.
  \item It discovers the BMS GATT service UUID \texttt{0000ffe0-...} and
        subscribes to two characteristics: FFE1 for telemetry notifications
        and FFE2 for writing control commands.
\end{enumerate}

The JK BMS does not transmit telemetry in a single BLE packet. It sends a
300-byte payload fragmented across multiple 20-byte BLE notification packets,
which is the maximum payload size permitted per BLE notification. The
\texttt{notifyCallback} function is registered to fire on every incoming
notification and reassembles these fragments as follows:

\begin{lstlisting}[
    language=CppESP,
    caption={BLE Frame Reassembly Logic}
]
void notifyCallback(BLERemoteCharacteristic* pChar,
                    uint8_t* pData, size_t length, bool isNotify) {
    if (length >= 4 &&
        pData[0] == 0x55 && pData[1] == 0xAA &&
        pData[2] == 0xEB && pData[3] == 0x90) {
        frameLen = 0;
    }
    for (size_t i = 0; i < length && frameLen < 300; i++) {
        frameBuffer[frameLen++] = pData[i];
    }
    if (frameLen >= 300) {
        if (frameBuffer[4] == 0x02) {
            parseCellInfo(frameBuffer);
        }
        frameLen = 0;
    }
}
\end{lstlisting}

\subsubsection{JK BMS Binary Protocol Parser}

The \texttt{parseCellInfo()} function decodes the raw 300-byte binary frame
using hardcoded byte offsets from the JK BMS JK02\_32S protocol specification,
with a frame offset constant \texttt{O = 32}. Table~\ref{tab:bmsbytes} shows
the key byte mappings.

\begin{longtable}{L{4.2cm} L{5.2cm} C{2.5cm}}
\caption{JK BMS Frame Key Byte Offsets (JK02\_32S, O = 32)}
\label{tab:bmsbytes}\\

\toprule
\textbf{Parameter} & \textbf{Decoding Formula} & \textbf{Byte Offset} \\
\midrule
\endfirsthead

\multicolumn{3}{c}%
{{\tablename\ \thetable{} -- continued }} \\
\toprule
\textbf{Parameter} & \textbf{Decoding Formula} & \textbf{Byte Offset} \\
\midrule
\endhead

\midrule
\multicolumn{3}{r}{{}} \\
\endfoot

\bottomrule
\endlastfoot

Cell Voltage [$n$]     & \texttt{get16(d, 6+$n\times$2) $\times$ 0.001\,V}      & $6 + n\times 2$ \\
Cell Resistance [$n$]  & \texttt{get16(d, 80+$n\times$2) $\times$ 0.001\,$\Omega$} & $80 + n\times 2$ \\
MOS Temperature        & \texttt{getS16(d, 112+O) $\times$ 0.1\,$^{\circ}$C}    & 144 \\
Total Pack Voltage     & \texttt{get32(d, 118+O) $\times$ 0.001\,V}             & 150 \\
Pack Power             & \texttt{getS32(d, 122+O) $\times$ 0.001\,W}            & 154 \\
Pack Current           & \texttt{getS32(d, 126+O) $\times$ 0.001\,A}            & 158 \\
Temperature 1          & \texttt{getS16(d, 130+O) $\times$ 0.1\,$^{\circ}$C}    & 162 \\
Temperature 2          & \texttt{getS16(d, 132+O) $\times$ 0.1\,$^{\circ}$C}    & 164 \\
Error Bitmask          & \texttt{get32(d, 134+O)}                                & 166 \\
State of Charge        & \texttt{d[141+O]\,\%}                                   & 173 \\
Cycle Count            & \texttt{get32(d, 150+O)}                                & 182 \\
State of Health        & \texttt{d[158+O]\,\%}                                   & 190 \\
Charging MOSFET        & \texttt{d[166+O] != 0 (bool)}                           & 198 \\
Discharging MOSFET     & \texttt{d[167+O] != 0 (bool)}                           & 199 \\
Balancing Active       & \texttt{d[169+O] != 0 (bool)}                           & 201 \\

\end{longtable}

After parsing, a CRC validation is applied. The 8-bit sum of all 299 preceding
bytes must equal byte \texttt{d[299]}. If the CRC fails, the frame is discarded
and \texttt{bmsData.valid} is set to \texttt{false} to prevent stale data from
propagating downstream.

\subsubsection{UART Inter-Board Communication}

Valid parsed data is serialised into a compact JSON object using the
ArduinoJson library and transmitted over UART2 (GPIO 17 TX, 115200 baud)
every five seconds. A newline character serves as the frame delimiter, which
allows Board 2 to detect a complete JSON frame reliably using a
readline-style receive loop.

\begin{lstlisting}[
    language=JavaScript,
    caption={UART Telemetry JSON Schema (Board 1 to Board 2)},
    label={lst:uartjson}
]
{
  "soc": 85, "soh": 100,
  "v": "13.200",  "i": "12.450",  "p": "161.24",
  "cap": "86.50", "fcap": "100.00",
  "t1": "28.3",   "t2": "29.1",   "mos": "31.5",
  "bal": "0.012", "cyc": 4,
  "cm": false,    "dm": true,     "err": 0,
  "rssi": -62,    "dv": "0.003",
  "cv": [3.324, 3.325, 3.326, 3.327],
  "cr": [0.120, 0.118, 0.122, 0.119]
}
\end{lstlisting}

\subsubsection{TFT Display Rendering}

The 2.4-inch ILI9341 TFT display is driven by the TFT\_eSPI library, which
provides hardware-accelerated SPI graphics on the ESP32. The display renders
a four-page carousel updated every two seconds:

Page 1 shows the core battery metrics: State of Charge as a large percentage
with a colour-coded status bar (green above 60\%, yellow between 20\% and 60\%,
red below 20\%), pack voltage, pack current and power.

Page 2 shows the three temperature sensor readings (Temperature 1, Temperature 2
and MOSFET temperature) with indicators flagging when any reading approaches the
protection limit.

Page 3 shows the individual voltage of each of the four cells, with the
highest and lowest voltage cells highlighted to assist in spotting imbalance.

Page 4 shows load prioritization status: which output channels are currently
active, which have been shed and the current SoC reading that the shedding
engine is acting on.

\begin{figure}[H]
  \centering
  \begin{subfigure}[t]{0.48\textwidth}
    \centering
    \includegraphics[
      width=0.8\textwidth,
      angle=270
    ]{The splash screen of the tft module 2.jpg}
    \caption{System splash screen on startup}
    \label{fig:tft_splash}
  \end{subfigure}
  \hfill
  \begin{subfigure}[t]{0.48\textwidth}
    \centering
    \includegraphics[
      width=0.8\textwidth,
      angle=270
    ]{Tft screen showing the first page of the carousel which contains the SOC, voltage, power and current_.jpg}
    \caption{Carousel page 1: SoC, voltage, current and power}
    \label{fig:tft_p1}
  \end{subfigure}

  \vspace{0.3cm}

  \begin{subfigure}[b]{0.48\textwidth}
    \centering
    \includegraphics[
      width=0.8\textwidth,
      angle=270
    ]{The tft showing the second screen of the carousel showing the temperatures, min and max voltage cell.jpg}
    \caption{Carousel page 2: temperatures and cell extremes}
    \label{fig:tft_p2}
  \end{subfigure}

  \caption{TFT display carousel: splash screen and first two telemetry pages}
  \label{fig:tft_pages_1}
\end{figure}

\begin{figure}[H]
  \centering
  \begin{subfigure}[b]{0.4\textwidth}
    \centering
    \includegraphics[width=\textwidth, angle=270]{The tft screen showing individual voltage LS of the 4 cells.jpg}
    \caption{Carousel page 3: individual cell voltages}
    \label{fig:tft_p3}
  \end{subfigure}
  \hfill
  \begin{subfigure}[b]{0.4\textwidth}
    \centering
    \includegraphics[width=\textwidth, angle=270, trim= 0cm 0cm 0cm 10cm, clip]{The tft showing the third carousel showing the load prioritization_.jpg}
    \caption{Carousel page 4: load prioritization status}
    \label{fig:tft_p4}
  \end{subfigure}
  \caption{TFT display carousel pages 3 and 4}
  \label{fig:tft_pages_2}
\end{figure}

\vspace{1cm}

\subsection{ESP32 Board 2 Firmware}

\subsubsection{UART Reception and JSON Parsing}

Board 2 listens on UART2 (GPIO 16 RX, 115200 baud) for newline-delimited
JSON frames from Board 1. When a complete frame arrives, the ArduinoJson
library deserialises it and populates the \texttt{BmsData} structure used
by the load shedding engine and the cloud uploader. Frames that fail JSON
parsing or do not match the expected schema are discarded without affecting
the relay states.

\subsubsection{Intelligent Load Shedding Engine}

The load shedding engine implements sequential priority-based load
disconnection with hysteresis. The three output channels are assigned the
priorities shown in Table~\ref{tab:priorities}.

\begin{table}[H]
\centering
\caption{Load Channel Priority Assignments and SoC Thresholds}
\label{tab:priorities}
\begin{tabular}{C{1.5cm} L{3.5cm} C{3.0cm} C{3.0cm}}
\toprule
\textbf{Channel} & \textbf{Load Type} & \textbf{Shed Threshold} & \textbf{Restore Threshold} \\
\midrule
CH3 & Non-essential (AC, entertainment) & SoC $<$ 60\% & SoC $>$ 65\% \\
CH2 & Major (fans, TVs, computers) & SoC $<$ 40\% & SoC $>$ 45\% \\
CH1 & Critical (routers, lighting, medical) & SoC $<$ 20\% & SoC $>$ 25\% \\
\bottomrule
\end{tabular}
\end{table}

The five-percentage-point hysteresis band between the shed and restore
thresholds prevents rapid relay oscillation when the SoC hovers near a
threshold boundary. A channel that has been shed stays off until the SoC
genuinely recovers five points above the shed level, confirming that the
battery is actually regaining charge before the load is reconnected.

\begin{lstlisting}[language=CppESP, caption={Load Shedding State Machine}]
void updateLoadShedding(int soc) {
    if (loadManager.mode != AUTO) return;
    if (soc < SOC_SHED_CH3 && ch3Active)
        { setRelay(CH3, OFF); ch3Active = false; }
    else if (soc > SOC_RESTORE_CH3 && !ch3Active)
        { setRelay(CH3, ON);  ch3Active = true;  }
    if (soc < SOC_SHED_CH2 && ch2Active)
        { setRelay(CH2, OFF); ch2Active = false; }
    else if (soc > SOC_RESTORE_CH2 && !ch2Active && ch3Active)
        { setRelay(CH2, ON);  ch2Active = true;  }
    if (soc < SOC_SHED_CH1 && ch1Active)
        { setRelay(CH1, OFF); ch1Active = false; }
    else if (soc > SOC_RESTORE_CH1 && !ch1Active)
        { setRelay(CH1, ON);  ch1Active = true;  }
}
\end{lstlisting}

This load shedding engine directly implements Objective~(iii).
\vspace{1cm}
\subsubsection{ACS758 Current Measurement}

\begin{lstlisting}[language=CppESP, caption={ACS758 Current Calculation}, label={lst:acs758}]
float readChannelCurrent(int adcPin, float zeroRef) {
    long sum = 0;
    for (int i = 0; i < 64; i++) {
        sum += analogRead(adcPin);
        delayMicroseconds(100);
    }
    float avgADC = sum / 64.0;
    float voltage = (avgADC / 4095.0) * 3300.0;  // mV
    float current = (voltage - zeroRef) / 40.0;   // 40 mV/A
    if (abs(current) < 0.10) current = 0.0;       // deadband
    return current;
}
\end{lstlisting}

\subsubsection{Firebase Realtime Database Synchronisation}

Board 2 maintains a Wi-Fi connection with automatic reconnection on dropout.
After connecting, it queries an NTP server (primary: \texttt{pool.ntp.org},
fallback: \texttt{time.google.com}, timezone UTC+1 for West Africa Time) to
synchronise its real-time clock. This timestamp feeds both the Supabase log
and the clock display forwarded to the TFT via the UART heartbeat.

Every five seconds the board publishes the full system state to Firebase using
an HTTPS PUT request to the Realtime Database REST endpoint. Board 2 also polls
the Firebase \texttt{/commands.json} and \texttt{/loadManager.json} paths every
three seconds using HTTPS GET requests. Commands written by the web application
to these paths are detected on the next poll cycle and executed within three
seconds.

\subsubsection{Supabase Historical Logging}

Every five seconds Board 2 also posts a timestamped row to the Supabase
\texttt{telemetry\_history} PostgreSQL table via HTTPS POST to the PostgREST
endpoint. NTP-synchronised Unix timestamps in milliseconds anchor every record
accurately in time. The stored columns cover voltage, current, power, all
three temperatures, all four cell voltages and the three per-channel load
currents. This historical table is what feeds the analytics charts in the
web application.

\subsection{Progressive Web Application}

\subsubsection{Overview and Technology Stack}

The Progressive Web Application is the user-facing interface of the entire
system. It is built as a Next.js 14 application written in TypeScript,
rendered as a client-side single-page application. The application is a
true PWA: it has a registered service worker and can be installed to the
home screen of any Android or iOS device via the browser's ``Add to Home
Screen'' prompt, giving it the look and feel of a native app without requiring
any app store submission.

The application communicates exclusively with the two cloud backends. It has
no direct connection to either ESP32. Firebase Realtime Database provides the
live telemetry stream and Supabase PostgreSQL provides the historical records
for analytics charts.

\subsubsection{Authentication}

Before any battery data is displayed, the user must sign in through Google
OAuth managed by Supabase Auth. On first visit, a login screen is shown with
a ``Continue with Google'' button. After successful authentication, Google
returns the user object to Supabase, which issues a JWT session token stored
client-side. A specific check for the system owner's email address unlocks
additional configuration controls not available to other viewers.

\subsubsection{Real-Time Data Flow}

On application load, a single Firebase \texttt{onValue} listener is registered
to the root of the Realtime Database. This listener fires immediately on every
write from ESP32 Board 2. Firebase pushes the update to the browser over a
persistent WebSocket connection, so no polling is required. The application
parses the incoming JSON into React state variables that drive the dashboard
display in real time.

A software watchdog monitors the \texttt{lastUpdate} counter that Board 2
increments on every telemetry upload. If no new value is received for more
than 30 seconds, the application declares the system offline and updates all
status indicators accordingly.

In parallel, the application fetches the last 500 historical records from
Supabase on startup and accumulates new snapshots every five seconds in memory.
This keeps the analytics charts growing with fresh data for as long as the
application is open.

\subsubsection{Navigation Structure}

The application uses a fixed bottom navigation bar with four tabs:

\begin{itemize}[noitemsep, leftmargin=0.7in]
  \item \textbf{Battery tab:} core battery telemetry including the SoC arc
        gauge, individual cell voltages, temperature readings and cycle count.
  \item \textbf{Energy tab:} real-time power flow diagram, remote load channel
        controls and runtime prediction.
  \item \textbf{Diagnostics tab:} active alarm panel and chronological system
        event log.
  \item \textbf{Profile tab:} user account information and application settings.
\end{itemize}

\subsubsection{Battery Tab and Real-Time Dashboard}

The Battery tab serves as the primary real-time monitoring interface for the user. Its central visual element is an animated, semicircular State of Charge (SoC) arc gauge that sweeps dynamically from $0^{\circ}$ to $180^{\circ}$ in proportion to remaining pack capacity, with color transitions indicating battery health (green above $60\%$, amber between $20\%$ and $60\%$, and red below $20\%$). Below the arc gauge, four dedicated telemetry metric tiles display pack terminal voltage, current, power, and remaining capacity, each providing an expand trigger to view historical time-series analytics. Individual cell voltages for all four series cells are rendered as comparative horizontal progress bars bounded between $2.70\,\text{V}$ and $3.60\,\text{V}$, prominently displaying the real-time cell voltage differential ($\text{DV}$) as an immediate indicator of pack balance. Figure~\ref{fig:app_screens} illustrates the three core user interfaces of the application: the Battery, Energy, and Diagnostics dashboards.

\begin{figure}[H]
  \centering

  \begin{subfigure}[b]{0.31\textwidth}
    \centering
    \includegraphics[width=0.75\textwidth]{batterytab.png}
    \caption{Battery tab}
    \label{fig:app_battery}
  \end{subfigure}
  \hfill
  \begin{subfigure}[b]{0.31\textwidth}
    \centering
    \includegraphics[width=0.75\textwidth]{Screenshot_20260828-144806.png}
    \caption{Energy tab}
    \label{fig:app_energy}
  \end{subfigure}
  \hfill
  \begin{subfigure}[b]{0.31\textwidth}
    \centering
    \includegraphics[width=0.75\textwidth]{The diagnostics tab page showing the activity feed and alarms feed.png}
    \caption{Diagnostics tab}
    \label{fig:app_diag}
  \end{subfigure}

  \caption{Web application screens: Battery, Energy and Diagnostics tabs}
  \label{fig:app_screens}
\end{figure}

\subsubsection{Remote Load Management}

The Energy tab contains the load management control panel. It has a mode
toggle that switches between AUTO and MANUAL operation. In AUTO mode, the
relay states reflect what Board 2 is doing autonomously based on the SoC
thresholds in the load shedding engine, and the channel controls are read-only.
In MANUAL mode, each channel gets an active toggle switch the user can operate
remotely.

When the user flips a channel toggle, the application writes immediately to
Firebase at \texttt{/loadManager}:

\begin{lstlisting}[language=JavaScript, caption={Load Manager Firebase Write on Manual Toggle}]
{ "mode": "MANUAL", "load1": 1, "load2": 0, "load3": 1 }
\end{lstlisting}

Board 2 reads this within its next three-second poll cycle and operates the
physical relays. The updated relay state propagates back to the web application
in the next five-second Firebase PUT, completing the control loop. A
twelve-second debounce prevents the Firebase listener from immediately
overwriting the user's manual selection with stale data during the period
between toggle and hardware acknowledgement.

\subsubsection{Emergency BMS Discharge Control}

The Energy tab also contains a BMS discharge MOSFET control card. This provides
a software-accessible emergency disconnect for the battery pack, independent of
the relay switching mechanism. When the user presses the disconnect button, the
application writes to two Firebase paths simultaneously:

\begin{lstlisting}[language=JavaScript, caption={BMS Discharge Command Write}]
// Sent to /commands
{ "discharge": false }
// Sent to /loadManager
{ "inverterOff": true }
\end{lstlisting}

Board 2 reads the \texttt{/commands} path within three seconds and forwards
the command as a JSON frame to Board 1 over UART. Board 1 composes the
appropriate JK BMS BLE command frame and writes it to the FFE2 GATT
characteristic, physically opening the discharge MOSFET. A Syncing badge
appears on the card while the commanded state differs from the confirmed
hardware state, indicating the command is still in transit through the UART
chain.

\subsubsection{Diagnostics: Alarms and Activity Feed}
The Diagnostics tab implements a two-tier fault monitoring architecture.
The first tier consists of four \emph{application-level threshold alarms}
computed directly from the live telemetry values received via Firebase.
The second tier exposes the JK BMS \emph{hardware error register} - a
32-bit bitmask written by the BMS firmware into its data frame and decoded
by ESP32 Board~1 during parsing.
\paragraph{Tier 1 -- Application Threshold Alarms.}
The four application-level alarms are evaluated client-side in the \texttt{ActiveAlarms}
component using the live state variables populated by the Firebase \texttt{onValue}
listener. They are re-evaluated on every telemetry update (approximately every
5~seconds). Table~\ref{tab:app_alarms} defines these conditions.
\begin{table}[H]
\centering
\caption{Tier~1: Application-Level Alarm Conditions}
\label{tab:app_alarms}
\begin{tabular}{L{3.8cm} L{5.2cm} C{2.5cm}}
\toprule
\textbf{Alarm Name} & \textbf{Trigger Condition} & \textbf{Severity} \\
\midrule
High Temperature   & Temperature sensor 1 $> 45\,^{\circ}$C           & Critical \\
Cell Imbalance     & $V_{\max} - V_{\min} > 15\,\text{mV}$             & Warning  \\
BMS Register Error & Hardware error bitmask $\neq$ 0                   & Critical \\
Low State of Charge & SoC $< 10\%$                                     & Warning  \\
\bottomrule
\end{tabular}
\end{table}
The cell delta is computed as:
\begin{equation}
    \Delta V_{\text{cell}} = \max(V_1, V_2, V_3, V_4) - \min(V_1, V_2, V_3, V_4)
    \label{eq:cell_delta}
\end{equation}
where $V_1$ through $V_4$ are the individual cell voltages received from
the BMS. A delta exceeding 15\,mV indicates a cell imbalance condition that
may reduce effective pack capacity if left unaddressed.
\paragraph{Tier 2 -- JK BMS Hardware Error Register.}
The JK BMS maintains a 32-bit error register that it encodes into bytes
166--169 of the JK02\_32S data frame. Each set bit corresponds to a
specific hardware fault condition. ESP32 Board~1 extracts this value
with a 32-bit little-endian read during frame parsing and forwards it
to the cloud as the \texttt{battery.errors.raw} field in Firebase.
The \texttt{ActiveAlarms} component in the PWA decodes this bitmask
bit-by-bit using the same 29-entry \texttt{errorNames} lookup table
that exists verbatim in the ESP32 firmware, ensuring perfect naming
consistency between the hardware fault register and the user interface.
Table~\ref{tab:bms_errors} lists all named fault bits.
\begin{longtable}{C{1.2cm} L{5.0cm} C{2.0cm} L{4.5cm}}
\caption{Tier~2: JK BMS 32-Bit Hardware Error Register  Named Fault Bits}
\label{tab:bms_errors} \\
\toprule
\textbf{Bit} & \textbf{Fault Name} & \textbf{Severity} & \textbf{Description} \\
\midrule
\endfirsthead
\multicolumn{4}{c}{\textit{(Table \ref{tab:bms_errors} continued)}} \\
\toprule
\textbf{Bit} & \textbf{Fault Name} & \textbf{Severity} & \textbf{Description} \\
\midrule
\endhead
\midrule
\multicolumn{4}{r}{\textit{}} \\
\endfoot
\bottomrule
\endlastfoot
0  & Wire resistance          & Warning  & Cell wire or contact resistance out of tolerance \\
1  & MOSFET overtemperature   & Warning  & Charge/discharge MOSFET die temperature too high \\
2  & Cell count mismatch      & Warning  & Detected cell count differs from configured count \\
3  & \textit{(Reserved)}      & ---      & Reserved bit; not decoded \\
4  & Battery fully charged    & Info     & Pack has reached full charge; informational flag \\
5  & Pack overvoltage         & Critical & Total pack voltage exceeded upper protection limit \\
6  & Charge overcurrent       & Warning  & Charge current exceeded the configured limit \\
7  & Charge short circuit     & Critical & Short circuit detected on the charge path \\
8  & Charge overtemperature   & Warning  & Battery temperature too high during charging \\
9  & Charge undertemperature  & Warning  & Battery temperature too low during charging \\
10 & Coprocessor comm error   & Warning  & Internal communication fault with BMS coprocessor \\
11 & Cell undervoltage        & Critical & At least one cell has fallen below minimum voltage \\
12 & Pack undervoltage        & Critical & Total pack voltage has dropped below protection limit \\
13 & Discharge overcurrent    & Warning  & Discharge current exceeded the configured limit \\
14 & Discharge short circuit  & Critical & Short circuit detected on the discharge path \\
15 & Discharge overtemperature& Warning  & Battery temperature too high during discharge \\
16 & Charging MOSFET abnormal & Warning  & Charge MOSFET failed to open or close correctly \\
17 & Discharging MOSFET abnormal & Warning & Discharge MOSFET failed to open or close correctly \\
18 & GPS disconnected         & Warning  & GPS module communication lost (JK BMS extended feature) \\
19 & Modify password in time  & Warning  & BMS security prompt to change default password \\
20 & Discharge on failed      & Warning  & Discharge enable command was not acknowledged \\
21 & Battery overtemperature  & Critical & Pack temperature has exceeded the absolute safety limit \\
22 & Temp sensor anomaly      & Critical & NTC temperature sensor reading out of valid range \\
23 & PL module anomaly        & Warning  & Auxiliary PL communication module fault \\
24 & SCP release failed       & Warning  & Short-circuit protection latch could not be released \\
25 & Discharge OCP II         & Critical & Secondary over-current protection level II triggered \\
26 & Discharge OCP III        & Critical & Tertiary over-current protection level III triggered \\
27 & Discharge undertemp alarm& Warning  & Discharge attempted below minimum temperature limit \\
28 & GPS remote lock          & Warning  & BMS remotely locked via GPS-linked command \\
29 & \textit{(Reserved)}      & ---      & Reserved bit; not decoded \\
30 & \textit{(Reserved)}      & ---      & Reserved bit; not decoded \\
31 & \textit{(Reserved)}      & ---      & Reserved bit; not decoded \\
\end{longtable}
Severity classifications are applied by the \texttt{getAlarmSeverity()}
function in \texttt{ActiveAlarms.tsx}. Faults classified as \emph{Critical}
indicate an imminent risk of damage or protection shutdown and are
highlighted in red. \emph{Warning} faults require attention but are not
immediately dangerous and are highlighted in amber. Bit~4
(Battery fully charged) is classified as \emph{Info} and renders in blue
as it is an expected operational state rather than a fault.
When the error bitmask is zero, the Diagnostics panel renders a green
``System Nominal'' state with a shield icon. When any bit is set, the
panel expands to show each active fault with its decoded name, severity
badge, a plain-language description and a recommended corrective action.
Tapping the expand button opens a full-screen modal listing all 32 bit
positions  active faults highlighted, clear faults dimmed  providing
a complete hardware register view for field diagnostics.
\paragraph{Voice Alert System.}
The application optionally reads active Tier~1 alarm descriptions aloud
using the Web Speech Synthesis API (\texttt{window.speechSynthesis}).
When enabled via the user preference \texttt{sentry\_audio\_warnings}
stored in \texttt{localStorage}, any non-zero \texttt{activeAlarmsCount}
triggers a spoken announcement, for example:
\textit{``Attention. Sentry has detected active telemetry alerts: high
temperature alarm, and active BMS register error.''} This feature
supports monitoring scenarios where the user cannot keep the display
in view at all times.
\paragraph{Activity Feed.}
Below the alarm panel, the Activity Feed renders a chronological event
log of all system state changes. Events are generated programmatically
by the application and fall into five categories:
\begin{enumerate}
\item \textbf{Boot events}, generated on application load when the
Firebase listener receives its first telemetry update, producing three
staggered log entries (400,ms apart): \textit{JK BMS BLE Connection
Active}, \textit{Sentry Gateway Online}, and \textit{Firebase Sync
Connected}.

\item \textbf{Connectivity events}, generated when the online
watchdog detects a transition to or from the offline state:
\textit{Sentry Gateway Offline} and \textit{BMS BLE Link Lost}.

\item \textbf{Relay control events}, generated when the user
manually toggles a load channel in the Smart Energy Manager, logging
which channel was switched and in which direction.

\item \textbf{Mode change events}, generated when the user switches
the load manager between \textsc{auto} and \textsc{manual} operating
modes.

\item \textbf{MOSFET command events}, generated when the user
activates or deactivates the BMS discharge MOSFET via the BMS
Discharge Control card.


\end{enumerate}

Each event record carries a severity badge (Success, Warning, Critical,
or Info), a category label (System, Hardware, Network, Manual, or BMS
Control), a title, a human-readable message and a relative timestamp
(e.g., \textit{``3\,min ago''}). Events are simultaneously forwarded to
the TopBar notification bell, where unread events accumulate as a count
badge until explicitly cleared.
\subsubsection{Historical Analytics Charts}

To empower the user with long-term performance observability and diagnostic tracking, tapping on any metric card on the Battery tab opens a dedicated, full-screen history chart modal. The system provides five specialized historical modals covering total pack voltage, current draw, power demand, thermal profile, and individual cell voltages. Each modal queries the synchronized \texttt{historyRecords} state array, which is hydrated on application startup from Supabase and updated continuously at 5-second intervals. The interface allows the user to toggle between a 24-hour perspective (120 evenly downsampled points) and a 7-day retrospective (168 hourly aggregated points), rendered as SVG spline curves with dynamic Y-axis scaling. Figure~\ref{fig:hist_charts} displays the historical modal interfaces for pack voltage and individual cell voltage tracking.

\begin{figure}[H]
  \centering
  \begin{subfigure}[b]{0.4\textwidth}
    \centering
    \includegraphics[width=\textwidth]{A card modal showing the pack voltage history_.png}
    \caption{Pack voltage history chart modal showing trend over the selected time window}
    \label{fig:hist_voltage}
  \end{subfigure}
  \hfill
  \begin{subfigure}[b]{0.4\textwidth}
    \centering
    \includegraphics[width=\textwidth]{A card modal showing each cell voltages History_.png}
    \caption{Cell voltages history chart modal showing individual cell trends}
    \label{fig:hist_cells}
  \end{subfigure}
  \caption{Historical analytics chart modals accessible from the Battery tab}
  \label{fig:hist_charts}
\end{figure}

The cell voltage chart uses fixed $2.70\,\text{V}$ to $3.60\,\text{V}$ scaling to prevent small millivolt-level noise from appearing exaggerated. Furthermore, each modal provides dedicated export utilities enabling the user to download the filtered time-series dataset as structured JSON or spreadsheet-compatible CSV files for offline academic and diagnostic analysis.

\section{System Integration and Testing}

\subsection{Controller Unit Assembly}

The physical controller unit centralizes all sensing, computation, and actuation electronics into a unified, modular assembly. Both ESP32 development boards, the 4-channel power relay module, the three ACS758 current sensors, and the DC-DC buck converter were integrated onto a Vero board substrate and enclosed within a robust, flame-retardant plastic enclosure. Interconnecting jumper wires and high-current power cables were routed along distinct physical paths to prevent electromagnetic coupling between high-current switching paths and sensitive analog sensing lines. Figure~\ref{fig:controller_assembly} shows the internal layout and wiring of the controller unit, while Figure~\ref{fig:controller_complete} illustrates the finalized controller package.

\vspace{1cm}
\begin{figure}[H]
  \centering

  \begin{subfigure}[t]{0.44\textwidth}
    \centering
    \includegraphics[
      height=6cm,
      angle=90,
      trim=0cm 0cm 2cm 0cm,
      clip
    ]{controller.jpg}
    \caption{Controller unit showing relays, dual ESP32 boards, current sensors and buck converter on the Vero board}
    \label{fig:controller_overview}
  \end{subfigure}
  \hspace{0.10\textwidth}
  \begin{subfigure}[t]{0.44\textwidth}
    \centering
    \includegraphics[
      height=6cm,
      angle=90
    ]{A pic of the whole component sitting in the enclosure with jumper wires connection 1.jpg}
    \caption{All controller components wired and seated inside the plastic enclosure}
    \label{fig:controller_wired}
  \end{subfigure}

  \caption{Controller unit assembly inside the project enclosure}
  \label{fig:controller_assembly}
\end{figure}

\begin{figure}[H]
  \centering
  \includegraphics[width=0.45\textwidth, angle=90]{A pic of the whole component sitting in the enclosure 2.jpg}
  \caption{Completed controller unit housed in the project enclosure, showing the relay stack, ESP32 boards and sensor wiring}
  \label{fig:controller_complete}
\end{figure}

\subsection{Integration Sequence}

System integration was carried out in stages to isolate problems at each
interface before adding the next layer. The battery pack was assembled and
tested independently first. The BMS was then connected and its Bluetooth
broadcast verified using the manufacturer's mobile application before any
ESP32 firmware was involved. Board 1 firmware was developed and tested
against the live BMS Bluetooth broadcast, confirming correct frame reassembly
and data parsing. UART communication between the two boards was verified using
serial monitor logging on both sides before cloud connectivity was added.
Board 2 Wi-Fi, Firebase and Supabase connectivity was tested with hardcoded
test payloads before the live UART-received data was introduced. The web
application was tested against live cloud data before relay control and MOSFET
command pathways were verified end-to-end.

\subsection{Test Scenarios}

The following test scenarios were defined and executed to verify system
behaviour:

\begin{enumerate}[noitemsep, leftmargin=0.7in]
  \item \textbf{BLE frame integrity:} Monitor BLE RSSI and the
        \texttt{bmsData.valid} flag over a two-hour continuous operation
        period to verify reliable frame reception and CRC validation.
  \item \textbf{UART data fidelity:} Compare values received by Board 2
        over UART against values read by Board 1 from the BMS to verify
        correct serialisation, transmission and deserialisation.
  \item \textbf{Current sensor calibration:} Apply known reference loads to
        each channel and compare ACS758 readings against a calibrated clamp
        meter.
  \item \textbf{Load shedding sequence:} Inject simulated SoC values into
        the shedding engine and verify that channels disconnect and reconnect
        at the correct thresholds in the correct order.
  \item \textbf{Remote command latency:} Measure the elapsed time between a
        relay toggle command issued from the web application and the physical
        relay actuation, repeated ten times to establish mean and maximum
        latency.
  \item \textbf{Cloud synchronisation continuity:} Monitor Firebase and
        Supabase write success rates over a 24-hour continuous operation
        period to verify reliable cloud publishing.
\end{enumerate}


\chapter{RESULTS AND DISCUSSION}

\section{Preamble}

This chapter presents the results obtained from the implementation and
testing of the intelligent inverter monitoring system. The results are
organized in alignment with the project objectives and the methodology
described in Chapter 3. Each result is discussed in context to explain
its significance, identify any deviations from expected behaviour and
assess the implications for overall system performance.

\section{Battery Pack Assembly and BMS Integration}

\subsection{Battery Pack Assembly Results}

The four 100Ah LiFePO\textsubscript{4} cells were tested individually
before assembly. All four cells registered voltages within 5\,mV of each
other, confirming acceptable uniformity for direct series connection without
requiring individual pre-charging. After top-balancing and series
connection, the assembled pack registered a total voltage of 13.16\,V on a
calibrated digital multimeter, consistent with the expected value for cells
at the tested State of Charge level.

The JK BMS power-on self-test completed successfully. Cell voltages
reported by the BMS over Bluetooth matched the individual cell voltages
measured by direct multimeter contact at each inter-cell junction to within
5\,mV, confirming accurate cell voltage sensing. The over-voltage and
under-voltage protection thresholds were verified by manually driving a
single cell toward its limits while monitoring the BMS MOSFET states in
the commercial BMS application. The charge MOSFET opened correctly at the
3.65\,V per cell overcharge threshold and the discharge MOSFET opened
correctly at the 2.5\,V per cell over-discharge threshold.

The active balancer operation was confirmed by allowing the pack to sit
at a partially charged state for two hours. The initial cell voltage
differential of 8\,mV reduced to 2\,mV over the balancing period,
demonstrating effective charge transfer from the higher-voltage cells to
the lower-voltage cells.

\subsection{Physical Construction and Electrical Interconnection Results}

The metallic enclosure housed the four cells with 3D-printed spacers
maintaining electrical isolation between adjacent cell casings. The
compression plates, tightened via threaded steel rods, applied uniform
pressure across the cell faces. No visible swelling was observed in any
cell after multiple charge and discharge cycles, indicating that the
mechanical constraint was effective.

All main power connections were made with 16\,mm\textsuperscript{2} cable
and crimped lugs at every terminal point. During a charging session with
a measured current of 26.75\,A, no abnormal heating was detected at any
joint by touch, confirming that the crimped lug connections presented
acceptably low contact resistance.

\subsection{BMS Configuration and Calibration Results}

The BMS protection parameters were configured as documented in
Table~\ref{tab:bmsconfig}. The calibration charge was carried out under
supervision until the tail current tapered to approximately 2\,A. At this
point the BMS application indicated 100\% State of Charge, confirming
that the coulomb counter reference had been successfully reset. The
configured COVP of 3.55\,V per cell and CUVP of 3.00\,V per cell both
lie within the safe operating window for LiFePO\textsubscript{4} and
provide deliberate margin against the chemistry's absolute limits.

\subsection{Inverter Coordination Results}

The Nexus 1.5\,kVA hybrid inverter was configured with a bulk charge
voltage of 14.5\,V and a charge current limit of 30\,A. During the
calibration charge, the BMS charge MOSFET remained closed until the pack
voltage approached the BMS's pack-level COVP equivalent, at which point
the BMS began limiting charge current before the inverter's 14.5\,V bulk
target was reached uniformly. This confirms the desired coordination
outcome in which the BMS acts as the more conservative safety limit.

\subsection{Discussion}

The small initial cell voltage differential of 8\,mV observed before
balancing is within the expected range for new prismatic cells from a
reputable manufacturer and does not indicate a quality concern. The
reduction to 2\,mV after two hours of active balancing confirms that the
JK BMS active balancer is functioning correctly and that the inter-cell
connections are providing adequate current paths for the energy transfer
circuit. The residual 2\,mV differential falls below the 5\,mV threshold
typically considered acceptable for well-balanced LiFePO\textsubscript{4}
packs.

The absence of heating at the crimped lug joints under 26.75\,A confirms
that the torque practice and cable sizing described in the methodology
were correctly implemented. The BMS calibration procedure successfully
established a full-charge reference for the coulomb counter, which is
essential for accurate State of Charge reporting in subsequent operation.

\section{Firmware Development and Communication Performance}

\subsection{BLE Communication and Frame Parsing Results}

Board 1 established a BLE connection to the JK BMS successfully on first
attempt in all ten trials conducted at varying distances from 1\,m to 8\,m.
The RSSI values ranged from $-42$\,dBm at 1\,m to $-75$\,dBm at 8\,m,
all within acceptable range for reliable BLE communication. No connection
failures occurred during a two-hour continuous monitoring session, and the
\texttt{bmsData.valid} flag remained \texttt{true} throughout, indicating
100\% CRC pass rate for received frames over the test period.

A representative parsed telemetry reading is shown in Table
\ref{tab:parseddata}.

\begin{longtable}{L{5.0cm} C{3.0cm} C{2.0cm}}
\caption{Representative Parsed JK BMS Telemetry Values}
\label{tab:parseddata}\\

\toprule
\textbf{Parameter} & \textbf{Parsed Value} & \textbf{Unit} \\
\midrule
\endfirsthead

\multicolumn{3}{c}%
{{\tablename\ \thetable{} -- continued }} \\
\toprule
\textbf{Parameter} & \textbf{Parsed Value} & \textbf{Unit} \\
\midrule
\endhead

\midrule
\multicolumn{3}{r}{{ }} \\
\endfoot

\bottomrule
\endlastfoot

Pack Voltage & 13.204 & V \\
Pack Current & 12.450 & A \\
Pack Power & 164.39 & W \\
State of Charge & 85 & \% \\
State of Health & 100 & \% \\
Temperature Sensor 1 & 28.3 & C \\
Temperature Sensor 2 & 29.1 & C \\
MOSFET Temperature & 31.5 & C \\
Cell Voltage Min & 3.298 & V \\
Cell Voltage Max & 3.302 & V \\
Voltage Differential & 4 & mV \\
Cycle Count & 4 & --- \\

\end{longtable}

\subsection{UART Inter-Board Communication Results}

Ten consecutive telemetry JSON frames transmitted from Board 1 to Board 2
were captured on both sides using serial monitor logging. All ten frames
were received intact and parsed successfully by Board 2. The numeric
values in the received JSON matched the values transmitted by Board 1 to
within floating-point rounding, confirming correct serialisation and
deserialisation. The five-second transmission interval was measured at an
average of 5.02 seconds over 20 consecutive transmissions, within 0.4\%
of the target.

\subsection{Discussion}

The 100\% CRC pass rate over the two-hour test session demonstrates that
the NimBLE-Arduino library and the custom frame reassembly logic handle
the fragmented 300-byte BLE payload reliably. The CRC validation is
essential: a corrupted frame that passes parsing would introduce incorrect
State of Charge values into the load shedding engine, potentially shedding
loads prematurely or failing to shed them when necessary.

The RSSI of $-75$\,dBm at 8\,m remains above the $-80$\,dBm threshold
typically cited for reliable BLE communication, suggesting the system
can operate reliably with the BMS and the ESP32 board separated by the
lengths typical of a household inverter installation.

The negligible timing deviation from the 5-second target is attributable
to the non-deterministic nature of the Arduino \texttt{millis()} timer
resolution and the variable time required for ArduinoJson serialisation.
This deviation is inconsequential for the application: the cloud
synchronisation interval is also five seconds, and a 20\,ms variation in
the UART transmission timing introduces no perceptible delay in the
dashboard update cycle.

\section{Intelligent Load Shedding and Current Sensing}

\subsection{ACS758 Current Sensor Accuracy Results}

Three known resistive loads were connected to Channel 1 in sequence and
the ACS758 readings were compared against a calibrated clamp meter.
Table~\ref{tab:sensoraccuracy} shows the results.

\begin{table}[H]
\centering
\caption{ACS758 Current Measurement Accuracy Comparison}
\label{tab:sensoraccuracy}
\begin{tabular}{C{3.0cm} C{3.0cm} C{3.0cm} C{2.5cm}}
\toprule
\textbf{Reference Load (A)} & \textbf{Clamp Meter (A)} & \textbf{ACS758 Reading (A)} & \textbf{Error (\%)} \\
\midrule
0.50 & 0.51 & 0.49 & 3.9 \\
1.20 & 1.22 & 1.19 & 2.5 \\
2.50 & 2.53 & 2.47 & 2.4 \\
\bottomrule
\end{tabular}
\end{table}

\subsection{Intelligent Load Shedding Performance Results}

The load shedding engine was tested by injecting simulated State of Charge
values at one-percent decrements from 100\% to 10\% and monitoring relay
actuation events. Table~\ref{tab:shedresults} shows the actual versus
expected actuation points.

\begin{table}[H]
\centering
\caption{Load Shedding and Restoration Actuation Verification}
\label{tab:shedresults}
\begin{tabular}{L{3.5cm} C{2.5cm} C{2.5cm} C{2.0cm} L{2.5cm}}
\toprule
\textbf{Event} & \textbf{Expected SoC} & \textbf{Actual SoC} & \textbf{Deviation} & \textbf{Result} \\
\midrule
CH3 shed & 60\% & 60\% & 0\% & PASS \\
CH2 shed & 40\% & 40\% & 0\% & PASS \\
CH1 shed & 20\% & 20\% & 0\% & PASS \\
CH1 restore & 25\% & 25\% & 0\% & PASS \\
CH2 restore & 45\% & 45\% & 0\% & PASS \\
CH3 restore & 65\% & 65\% & 0\% & PASS \\
\bottomrule
\end{tabular}
\end{table}

All six actuation events occurred at the correct thresholds with zero
deviation. The relay actuation sound (audible click) was confirmed for
each event. No spurious actuation events occurred between threshold
crossings during the test.

\subsection{Discussion}

The maximum measurement error of 3.9\% at the lowest test current is
within the acceptable range for load monitoring in this application. ACS758
sensors are specified at $\pm$1.5\% accuracy at full-scale current, but
at low currents the fixed offset error of the ADC contributes a larger
relative error. The 64-sample averaging implemented in the firmware
reduces random ADC noise effectively, as confirmed by the stable readings
with minimal fluctuation during each test measurement. For the purposes
of load management and dashboard display, the accuracy achieved is
sufficient.

The exact correspondence between expected and actual actuation thresholds
confirms that the integer State of Charge comparison logic in the shedding
engine is implemented correctly. The hysteresis band of five percentage
points between shed and restore thresholds proved effective in preventing
relay oscillation: during testing of the restore sequence, no channel was
re-shed after restoration until the SoC was explicitly driven back below
the shed threshold.

The relay click confirmed physical relay actuation as opposed to only a
logical state change in the firmware, which is an important distinction
for confidence in the hardware. In a live deployment, the ability to
precisely control at which State of Charge each load tier is shed gives
the user the confidence that critical loads will remain powered for as
long as the battery can sustain them.

\section{Cloud Synchronisation and Remote Control}

\subsection{Remote Command Latency Results}

Ten relay toggle commands were issued from the web application to Channel
2 (Major Load) and the time from command submission to audible relay
actuation was measured using a stopwatch. Results are presented in
Table~\ref{tab:latency}.

\begin{table}[H]
\centering
\caption{Remote Command Actuation Latency Measurements}
\label{tab:latency}
\begin{tabular}{C{2.0cm} C{3.0cm} C{2.0cm} C{3.0cm}}
\toprule
\textbf{Trial} & \textbf{Latency (s)} & \textbf{Trial} & \textbf{Latency (s)} \\
\midrule
1 & 2.1 & 6 & 1.9 \\
2 & 1.8 & 7 & 2.3 \\
3 & 2.4 & 8 & 1.7 \\
4 & 1.9 & 9 & 2.1 \\
5 & 2.2 & 10 & 2.0 \\
\midrule
\multicolumn{2}{l}{\textbf{Mean}} & \multicolumn{2}{l}{2.04 s} \\
\multicolumn{2}{l}{\textbf{Maximum}} & \multicolumn{2}{l}{2.4 s} \\
\multicolumn{2}{l}{\textbf{Minimum}} & \multicolumn{2}{l}{1.7 s} \\
\bottomrule
\end{tabular}
\end{table}

\subsection{Cloud Synchronisation Reliability Results}

Over a 24-hour continuous operation period, the system attempted 17,280
Firebase PUT operations (one every five seconds) and 17,280 Supabase POST
operations. Firebase recorded 17,256 successful writes, representing a
99.85\% success rate. Supabase recorded 17,241 successful inserts,
representing a 99.77\% success rate. All failures occurred in two clusters
corresponding to brief Wi-Fi disruptions caused by router restarts during
the test period.

\subsection{Discussion}

The mean command latency of 2.04 seconds is dominated by Board 2's
three-second Firebase polling interval. In the worst case, a command
written to Firebase arrives just after a poll cycle and must wait
approximately three seconds for the next poll. The measured maximum of
2.4 seconds is consistent with this interpretation, as the command write
to Firebase, the network round trip and the relay driver execution
together contribute an additional fraction of a second on top of the poll
interval remainder.

For a home inverter remote control application, a mean latency of
approximately two seconds is acceptable. The user experience of tapping
a button and seeing a relay respond within two seconds is responsive
enough for all practical load management scenarios. If sub-second latency
were required, replacing the polling mechanism with a Firebase real-time
listener or an MQTT subscription on Board 2 would reduce command latency
significantly at the cost of additional implementation complexity.

The 99.85\% and 99.77\% success rates demonstrate highly reliable cloud
synchronisation under normal operating conditions. The write failures
clustered around router restart events confirm that the automatic Wi-Fi
reconnection logic in the firmware is functioning: after the router
restarted, Board 2 re-established its Wi-Fi connection and resumed
publishing without any firmware reset or manual intervention. The brief
gap in the Supabase historical record during these disruptions is visible
in the analytics chart as a short flat segment, consistent with the
absence of data during the disconnected period.


%% ══════════════════════════════════════════════════════════════
%%  CHAPTER 5: CONCLUSION AND RECOMMENDATIONS
%% ══════════════════════════════════════════════════════════════
\newpage
\chapter{CONCLUSION AND RECOMMENDATIONS}

\section{Conclusion}

This project has successfully designed and implemented a complete
IoT-enabled intelligent inverter monitoring system demonstrating the
integration of LiFePO\textsubscript{4} battery technology, a JK Battery
Management System, a dual-ESP32 embedded firmware architecture and a
cloud-connected Progressive Web Application.

The dual-microcontroller architecture resolved a fundamental hardware
constraint of the ESP32 platform: the incompatibility between simultaneous
Bluetooth Low Energy and Wi-Fi operation on a single board. By dedicating
ESP32 Board 1 to BLE communication with the JK BMS and TFT display
rendering, and ESP32 Board 2 to Wi-Fi networking, relay switching,
current sensing and cloud publishing, the system achieves reliable
performance across all subsystems without radio interference.

The JK BMS binary protocol parser successfully decodes all telemetry
parameters from the proprietary 300-byte JK02\_32S frame with 100\% CRC
pass rate in testing. State of Charge, State of Health, cell voltages,
temperatures and MOSFET states are accurately extracted and transmitted
to the cloud. The intelligent load shedding engine correctly actuates
all six relay events (three sheds and three restores) at the precise
defined thresholds with hysteresis, demonstrating reliable priority-based
load management. The ACS758 current sensors provide per-channel current
measurements within 4\% of reference values, adequate for the load
monitoring application. Remote command latency averaged 2.04 seconds,
acceptable for household load control. Cloud synchronisation achieved
99.85\% write success rate over a 24-hour period.

Taken together, these results confirm that all four project objectives
have been achieved. An intelligent, remotely monitored inverter system
has been built and verified on physical hardware, demonstrating that such
a system is both technically achievable and practically deployable using
accessible and affordable hardware.

\section{Recommendations}

Based on the outcomes of this project, several avenues for future
improvement are recommended. First, implementing an MQTT subscription or
a Firebase Server-Sent Events listener on Board 2 would reduce remote
command latency from the current mean of approximately two seconds to well
below one second, significantly improving the responsiveness of the remote
control interface and eliminating the dependency on periodic polling.
Second, adding Firebase Cloud Messaging to the Progressive Web Application
would enable the system to alert the user via phone notification when
critical thresholds are crossed, for example when State of Charge drops
below 20\% or when battery temperature exceeds the warning threshold,
without requiring the web application to remain open in the browser.
Third, integrating an MPPT solar charge controller and monitoring its
output current through an additional ACS758 sensor would enable the web
application to display real-time solar generation alongside load
consumption, providing a complete energy balance picture that would
enhance the system's utility in off-grid and hybrid installations.
Finally, extending the battery pack capacity by adding cells in parallel
would extend runtime at full load and reduce the frequency of load
shedding events; the existing BMS and firmware architecture are designed
to accommodate a larger pack without significant modification, and
converting the Progressive Web Application into a native Android
application using React Native or Capacitor would provide better push
notification support, offline caching and a more polished user experience
than a browser-based PWA.

%% ══════════════════════════════════════════════════════════════
%%  REFERENCES
%% ══════════════════════════════════════════════════════════════

\newpage

\addcontentsline{toc}{chapter}{REFERENCES}

\setlength{\bibsep}{10pt}
\renewcommand{\bibname}{REFERENCES}

\begin{thebibliography}{99}
\setlength{\itemsep}{8pt}

\bibitem[\protect\citeauthoryear{Aliyu, Ramli and Saleh}{Aliyu et al.}{2013}]{aliyu2013}
Aliyu, A.~S., Ramli, A.~T., and Saleh, M.~A. (2013).
Nigeria electricity crisis: Power generation capacity expansion and
environmental ramifications.
\textit{Energy}, \textit{61}, 354--367.
https://doi.org/10.1016/j.energy.2013.09.021

\bibitem[\protect\citeauthoryear{Ashton}{Ashton}{2009}]{ashton2009}
Ashton, K. (2009).
That ``Internet of Things'' thing.
\textit{RFID Journal}, \textit{22}(7), 97--114.

\bibitem[\protect\citeauthoryear{Dini and Ria}{Dini and Ria}{2024}]{dini2024}
Dini, P., and Ria, A. (2024).
A compact overview on Li-Ion batteries characteristics and battery management
systems integration for automotive applications.
\textit{Energies}, \textit{17}(23), 5992.
https://doi.org/10.3390/en17235992

\bibitem[\protect\citeauthoryear{Divya and Ostergaard}{Divya and Ostergaard}{2009}]{divya2009}
Divya, K.~C., and Ostergaard, J. (2009).
Battery energy storage technology for power systems: An overview.
\textit{Electric Power Systems Research}, \textit{79}(4), 511--520.
https://doi.org/10.1016/j.epsr.2008.09.017

\bibitem[\protect\citeauthoryear{Dzahir and Chia}{Dzahir and Chia}{2023}]{dzahir2023}
Dzahir, M.~A.~S.~M., and Chia, K.~S. (2023).
Evaluating the energy consumption of ESP32 microcontroller for real-time
MQTT IoT-based monitoring system.
In \textit{2023 International Conference on Innovation and Intelligence for
Informatics, Computing, and Technologies (3ICT)} (pp. 255--261). IEEE.
https://doi.org/10.1109/3ICT60104.2023.10391358

\bibitem[\protect\citeauthoryear{Espressif Systems}{Espressif Systems}{2024}]{espressif2024}
Espressif Systems. (2024).
\textit{ESP32 technical reference manual}.
https://docs.espressif.com

\bibitem[\protect\citeauthoryear{Expo}{Expo}{2024}]{expodocs2024}
Expo. (2024).
\textit{Using Firebase with Expo}.
https://docs.expo.dev/guides/using-firebase/

\bibitem[\protect\citeauthoryear{Firebase}{Firebase}{2024}]{firebasedocs2024}
Firebase. (2024).
\textit{Cloud Firestore documentation}.
Google. https://firebase.google.com/docs/firestore

\bibitem[\protect\citeauthoryear{Fraunhofer}{Fraunhofer}{2022}]{fraunhofer2022}
Fraunhofer Institute for Solar Energy Systems. (2022).
\textit{Mechanical behaviour of lithium-ion prismatic cells under cycling
conditions}.
Fraunhofer ISE. https://publica.fraunhofer.de/entities/publication/a7531353-952b-4c94-8300-1a323943ab21

\bibitem[\protect\citeauthoryear{Hannan et al.}{Hannan et al.}{2017}]{hannan2017}
Hannan, M.~A., Hoque, M.~M., Hussain, A., Yusof, Y., and Ker, P.~J. (2017).
State-of-the-art and energy management system of lithium-ion batteries in
electric vehicle applications: Issues and recommendations.
\textit{IEEE Access}, \textit{6}, 19362--19378.
https://doi.org/10.1109/ACCESS.2018.2817655

\bibitem[\protect\citeauthoryear{IJCRT}{IJCRT}{2025}]{ijcrt2025}
IJCRT. (2025).
IoT-based battery monitoring system using BMS and ESP32 with Bluetooth
communication.
\textit{International Journal of Creative Research Thoughts}, \textit{13}(5).
https://www.ijcrt.org/papers/IJCRT2605567.pdf

\bibitem[\protect\citeauthoryear{IJPEDS}{IJPEDS}{2023}]{ijpeds2023}
International Journal of Power Electronics and Drive Systems. (2023).
State-of-charge-based intelligent load management for battery-backed
inverter systems.
\textit{IJPEDS}, \textit{14}(1).
https://ijpeds.iaescore.com/index.php/IJPEDS/article/view/23006

\bibitem[\protect\citeauthoryear{KIT}{KIT}{2024}]{kit2024}
Karlsruhe Institute of Technology. (2024).
\textit{Mechanical compression effects on LiFePO4 prismatic cell cycle life}.
KIT Repository.
https://publikationen.bibliothek.kit.edu/1000192706/180060941

\bibitem[\protect\citeauthoryear{Nitta et al.}{Nitta et al.}{2015}]{nitta2015}
Nitta, N., Wu, F., Lee, J.~T., and Yushin, G. (2015).
Li-ion battery materials: Present and future.
\textit{Materials Today}, \textit{18}(5), 252--264.
https://doi.org/10.1016/j.mattod.2014.10.040

\bibitem[\protect\citeauthoryear{OASIS}{OASIS}{2019}]{oasis2019}
OASIS. (2019).
\textit{MQTT Version 5.0}.
https://docs.oasis-open.org/mqtt/mqtt/v5.0/mqtt-v5.0.html

\bibitem[\protect\citeauthoryear{Omar et al.}{Omar et al.}{2014}]{omar2014}
Omar, N., Monem, M.~A., Firouz, Y., Salminen, J., Smekens, J., Hegazy, O.,
Gaulous, H., Mulder, G., Van den Bossche, P., Coosemans, T., and
Van Mierlo, J. (2014).
Lithium iron phosphate based battery: Assessment of the aging parameters and
development of cycle life model.
\textit{Applied Energy}, \textit{113}, 1575--1585.
https://doi.org/10.1016/j.apenergy.2013.09.003

\bibitem[\protect\citeauthoryear{Pillet et al.}{Pillet et al.}{2025}]{pillet2025}
Pillet, M., Zhang, L., and Haran, K. (2025).
Long-term cycling performance of LiFePO4 cells under variable temperature
and charge rate conditions.
\textit{Journal of the Electrochemical Society}, \textit{171}(7), 070510.
https://ui.adsabs.harvard.edu/abs/2024JElS..171g0510P/abstract

\bibitem[\protect\citeauthoryear{Piller, Perrin and Jossen}{Piller et al.}{2001}]{piller2001}
Piller, S., Perrin, M., and Jossen, A. (2001).
Methods for state-of-charge determination and their applications.
\textit{Journal of Power Sources}, \textit{96}(1), 113--120.
https://doi.org/10.1016/S0378-7753(01)00560-2

\bibitem[\protect\citeauthoryear{Rashid}{Rashid}{2014}]{rashid2014}
Rashid, M.~H. (2014).
\textit{Power electronics handbook: Devices, circuits, and applications}
(3rd ed.).
Butterworth-Heinemann.

\end{thebibliography}
\end{document}